% Options for packages loaded elsewhere
\PassOptionsToPackage{unicode}{hyperref}
\PassOptionsToPackage{hyphens}{url}
%
\documentclass[
  11pt,
]{article}
\usepackage{amsmath,amssymb}
\usepackage{iftex}
\ifPDFTeX
  \usepackage[T1]{fontenc}
  \usepackage[utf8]{inputenc}
  \usepackage{textcomp} % provide euro and other symbols
\else % if luatex or xetex
  \usepackage{unicode-math} % this also loads fontspec
  \defaultfontfeatures{Scale=MatchLowercase}
  \defaultfontfeatures[\rmfamily]{Ligatures=TeX,Scale=1}
\fi
\usepackage{lmodern}
\ifPDFTeX\else
  % xetex/luatex font selection
  \IfFontExistsTF{Times New Roman}{\setmainfont{Times New Roman}}{\setmainfont{DejaVu Serif}}
  \IfFontExistsTF{Arial}{\setsansfont{Arial}}{\setsansfont{DejaVu Sans}}
  \IfFontExistsTF{Courier New}{\setmonofont{Courier New}}{\setmonofont{DejaVu Sans Mono}}
\fi
% Use upquote if available, for straight quotes in verbatim environments
\IfFileExists{upquote.sty}{\usepackage{upquote}}{}
\IfFileExists{microtype.sty}{% use microtype if available
  \usepackage[]{microtype}
  \UseMicrotypeSet[protrusion]{basicmath} % disable protrusion for tt fonts
}{}
\makeatletter
\@ifundefined{KOMAClassName}{% if non-KOMA class
  \IfFileExists{parskip.sty}{%
    \usepackage{parskip}
  }{% else
    \setlength{\parindent}{0pt}
    \setlength{\parskip}{6pt plus 2pt minus 1pt}}
}{% if KOMA class
  \KOMAoptions{parskip=half}}
\makeatother
\usepackage{xcolor}
\usepackage[margin=20mm]{geometry}
\usepackage{longtable,booktabs,array}
\usepackage{calc} % for calculating minipage widths
% Correct order of tables after \paragraph or \subparagraph
\usepackage{etoolbox}
\makeatletter
\patchcmd\longtable{\par}{\if@noskipsec\mbox{}\fi\par}{}{}
\makeatother
% Allow footnotes in longtable head/foot
\IfFileExists{footnotehyper.sty}{\usepackage{footnotehyper}}{\usepackage{footnote}}
\makesavenoteenv{longtable}
\setlength{\emergencystretch}{3em} % prevent overfull lines
\providecommand{\tightlist}{%
  \setlength{\itemsep}{0pt}\setlength{\parskip}{0pt}}
\setcounter{secnumdepth}{5}
\ifLuaTeX
\usepackage[bidi=basic]{babel}
\else
\usepackage[bidi=default]{babel}
\fi
\babelprovide[main,import]{russian}
\ifPDFTeX
\else
\IfFontExistsTF{Times New Roman}{\babelfont{rm}{Times New Roman}}{\babelfont{rm}{DejaVu Serif}}
\fi
% get rid of language-specific shorthands (see #6817):
\let\LanguageShortHands\languageshorthands
\def\languageshorthands#1{}
% Исходник для редактирования в Codex. Кодировка UTF-8.
% Сборка: xelatex porous_literature_index.tex (запустить дважды).
% OCR и исходные номера страниц сохранены; сверять спорные места по сканам.
\usepackage{microtype}
\setlength{\emergencystretch}{3em}
\ifLuaTeX
  \usepackage{selnolig}  % disable illegal ligatures
\fi
\usepackage{bookmark}
\IfFileExists{xurl.sty}{\usepackage{xurl}}{} % add URL line breaks if available
\urlstyle{same}
\hypersetup{
  pdftitle={Указатель литературы по течениям в пористых средах},
  pdflang={ru},
  hidelinks,
  pdfcreator={LaTeX via pandoc}}

\title{Указатель литературы по течениям в пористых средах}
\author{}
\date{}

\begin{document}
\maketitle

{
\setcounter{tocdepth}{3}
\tableofcontents
}
Дата составления: 6 ноября 2026 года. Всего приложено 14 файлов: 12 книжных файлов, программа и рукописный конспект. Полные тексты книг не изучались; указатель составлен по титульным страницам, оглавлениям и программе.

% ===== Раздел 1 =====

\section{Как использовать}\label{section-01}

По вопросу найти подходящую книгу, главу и параграф; затем проверить нужное место по оригиналу. В ответе давать название книги, главу, параграф, печатную страницу и при необходимости страницу файла. Нумерация переводов и английских изданий различается. Извлеченные ниже оглавления содержат ошибки OCR, особенно в цифрах; спорные номера сверять по скану. Маркер «Страница файла» обозначает страницу PDF/DjVu, остальные номера внутри оглавлений --- печатные страницы.

% ===== Раздел 2 =====

\section{Список файлов}\label{section-02}

\begin{longtable}[]{@{}
  >{\raggedright\arraybackslash}p{(\columnwidth - 4\tabcolsep) * \real{0.3333}}
  >{\raggedright\arraybackslash}p{(\columnwidth - 4\tabcolsep) * \real{0.3333}}
  >{\raggedright\arraybackslash}p{(\columnwidth - 4\tabcolsep) * \real{0.3333}}@{}}
\toprule\noalign{}
\begin{minipage}[b]{\linewidth}\raggedright
№
\end{minipage} & \begin{minipage}[b]{\linewidth}\raggedright
Книга или материал
\end{minipage} & \begin{minipage}[b]{\linewidth}\raggedright
Файл
\end{minipage} \\
\midrule\noalign{}
\endhead
\bottomrule\noalign{}
\endlastfoot
1 & И. А. Чарный. Подземная гидрогазодинамика. 1963 & \texttt{Charny.djvu} \\
2 & А. Г. Куликовский, Е. И. Свешникова. Нелинейные волны в упругих телах. 1998 & \texttt{Kulikovskii-Sveshnikova.djvu} \\
3 & R. E. Collins. Flow of Fluids Through Porous Materials. Английское издание & \texttt{Collins-eng.pdf} \\
4 & Г. И. Баренблатт, В. М. Ентов, В. М. Рыжик. Теория нестационарной фильтрации жидкости и газа. 1972 & \texttt{Barenblatt-et-al-1972.djvu} \\
5 & К. С. Басниев, А. М. Власов, И. Н. Кочина, В. М. Максимов. Подземная гидравлика. 1986 & \texttt{Basniev-et-al.djvu} \\
6 & Н. Е. Леонтьев. Основы теории фильтрации. 2-е изд., 2017 & \texttt{Leontiev-2017.pdf} \\
7 & М. Маскет. Течение однородных жидкостей в пористой среде. 2004 & \texttt{Muskat-2004.djvu} \\
8 & Р. Коллинз. Течения жидкостей через пористые материалы. 1964 & \texttt{Collins.djvu} \\
9 & Myron B. Allen. The Mathematics of Fluid Flow Through Porous Media. 2021 & \texttt{Allen-The-mathematics-of-fluid-flow-through-porous.pdf} \\
10 & Jacob Bear. Dynamics of Fluids in Porous Media. 1972 & \texttt{Bear-Dynamics-of-fluids-in-porous-media-1972.djvu} \\
11 & Г. И. Баренблатт, В. М. Ентов, В. М. Рыжик. Движение жидкостей и газов в природных пластах. 1984 & \texttt{Barenblatt-et-al-1984.pdf} \\
12 & M. Muskat. The Flow of Homogeneous Fluids Through Porous Media. Английское издание & \texttt{Muskat-eng.djvu} \\
13 & Течения в пористых средах. Программа курса Н. Е. Леонтьева & \texttt{porous-program.pdf} \\
14 & Рукописный конспект лекций, 33 страницы PDF (развороты) & \texttt{porous-lecture-notes.pdf} \\
\end{longtable}

Важные различия: приложенный русский Маскет датирован 2004 годом (628 с.), хотя программа ссылается на издание 2006 года. Книга П. Я. Полубариновой-Кочиной «Теория движения грунтовых вод» (1977) есть в программе, но среди приложений отсутствует.

% ===== Раздел 3 =====

\section{Быстрый тематический поиск}\label{section-03}

\begin{longtable}[]{@{}
  >{\raggedright\arraybackslash}p{(\columnwidth - 2\tabcolsep) * \real{0.5000}}
  >{\raggedright\arraybackslash}p{(\columnwidth - 2\tabcolsep) * \real{0.5000}}@{}}
\toprule\noalign{}
\begin{minipage}[b]{\linewidth}\raggedright
Тема
\end{minipage} & \begin{minipage}[b]{\linewidth}\raggedright
Где начать поиск
\end{minipage} \\
\midrule\noalign{}
\endhead
\bottomrule\noalign{}
\endlastfoot
Пористость, скорость фильтрации, закон Дарси & Леонтьев, гл. I, §§ 1--4; Чарный, гл. I; Басниев, гл. 1; Allen, § 3.1 \\
Осреднение, балансы массы, импульса, энергии, граничные условия & Леонтьев, гл. I, §§ 3, 5, 6; Чарный, гл. II; Bear, гл. 4, 6, 7 \\
Анизотропия и нелинейные законы фильтрации & Леонтьев, §§ 3--4; Allen, §§ 3.2, 3.7; Bear, §§ 5.3, 5.6 \\
Скважины, стационарные течения, комплексный потенциал & Леонтьев, §§ 7--8; Чарный, гл. III--V; Маскет, гл. IV--V, IX; Allen, § 4.1 \\
Безнапорная фильтрация, Буссинеск, Дюпюи & Леонтьев, § 9; Чарный, гл. VI; Bear, гл. 8; Allen, § 4.3 \\
Фильтрация газа, уравнение Лейбензона & Леонтьев, § 11; Маскет, гл. XI; Басниев, гл. 5, 7 \\
Упругий режим, пьезопроводность, пуск скважины & Леонтьев, § 12; Чарный, гл. VIII; Баренблатт и др. (1972), гл. II--V \\
Неизотермическая фильтрация & Леонтьев, § 13; Баренблатт и др. (1972), гл. IX \\
Поршневое вытеснение, устойчивость фронта & Басниев, гл. 8, §§ 1, 2, 4; Чарный, гл. VII, §§ 6--7; Allen, § 6.5 \\
Фазовые проницаемости, капиллярное давление, функция Леверетта & Басниев, гл. 1, § 5 и гл. 9, §§ 2--3; Баренблатт и др. (1972), гл. VI, § 1; Чарный, гл. IX, § 1 \\
Бакли--Леверетт, скачок насыщенности & Басниев, гл. 9, §§ 4--5; Баренблатт и др. (1972), гл. VI, § 2; Чарный, гл. IX, § 2; Allen, § 6.4 \\
Эволюционность разрывов & Куликовский--Свешникова, §§ 1.3, 1.6 \\
Рапопорт--Лис, структура скачка насыщенности & Баренблатт и др. (1972), гл. VI, § 3; Басниев, гл. 9, § 6 \\
Трещиновато-пористые среды, двойная пористость & Баренблатт и др. (1972), гл. VII; Басниев, гл. 12; Баренблатт и др. (1984), гл. 3, § 4 \\
Перенос примеси, дисперсия, адсорбция & Allen, гл. 5; Bear, гл. 10; Коллинз, гл. 8; Баренблатт и др. (1984), гл. 5 \\
Численные методы, вариационная постановка & Allen, § 3.6; Bear, § 7.9; Басниев, гл. 13 \\
\end{longtable}

Ссылки на параграфы тем 9--12 сверены с программой курса; остальные строки --- навигация по оглавлениям, а не вывод о полноте изложения.

% ===== Раздел 4 =====

\section{Программа курса}\label{section-04}

{[}FILE PAGE 1{]} Течения в пористых средах Леонтьев Н. Е. 1. Основные физические сведения о пористых средах. Описание фильтрации в рамках механики сплошных сред. Пористость. Скорость фильтрации. 2. Закон Дарси и пределы его применимости. Учёт анизотропии пористой среды. Двучленный закон фильтрации Форхгеймера. Фильтрация с предельным градиентом давления. 3. Уравнения баланса массы, импульса и энергии для течений в пористых средах. Соотношения на разрывах. Типич- ные граничные условия. 4. Замкнутая система уравнений для фильтрации несжимаемой жидкости. Стационарные решения для одномерных течений (с плоскими волнами, течений в случае осевой симметрии). Напорный приток к совершенной скважине с постоянным расходом (дебитом). Формула Дюпюи для напорной фильтрации. 5. Плоская установившаяся фильтрация несжимаемой жидкости. Типичные постановки задач. Применение методов ТФКП. Примеры: постановка задачи о водонепроницаемой плотине с плоским основанием; постановка задачи о работе кротового оросителя. 6. Гидравлическая теория безнапорной фильтрации. Уравнение Буссинеска. Формула Дюпюи для безнапорной филь- трации. 7. Изотермическая фильтрация газа в недеформируемой пористой среде. Уравнение для давления (уравнение Лей- бензона). Аналогия с пологими безнапорными движениями несжимаемой жидкости в пористой среде. 8. Упругий режим фильтрации. Уравнение пьезопроводности. Постановка задачи о распространении плоской вол- ны повышения давления в полубесконечном пласте. Оценка характерного времени распространения возмущения давления. 9. Вытеснение несмешивающихся жидкостей. Одномерная задача о поршневом вытеснении несмешивающихся жид- костей (модель Лейбензона --- Маскета). Устойчивость фронта вытеснения (элементарная теория). {[}3, гл. 8, § 1, 2, 4; 8, гл. VII, § 6, 7{]} 10. Описание многофазной фильтрации. Обобщенный закон Дарси, относительная фазовая проницаемость. Учет ка- пиллярной разности давлений между фазами, функция Леверетта. {[}2, гл. VI, § 1; 3, гл. 1, § 5; гл. 9, § 2, 3; 8, гл. IX, § 1{]} 11. Модель двухфазного течения в пренебрежении капиллярной разностью давлений (модель Маскета --- Мереса). Опрокидывание волны повышения насыщенности в задаче об одномерном течении (в задаче Бакли --- Леверетта). Условия на скачках насыщенности. {[}2, гл. VI, § 2; 3, гл. 9, § 4, 5; 8, гл. IX, § 2{]} 12. Понятие об эволюционности разрывов для гиперболических уравнений. Эволюционность скачка насыщенности в задаче Бакли --- Леверетта. Модель двухфазных течений с учетом капиллярной разности давлений. Понятие об уравнении Рапопорта --- Лиса. Представление о задаче о структуре скачка насыщенности. {[}2, гл. VI, § 3; 3, гл. 9, § 6; 10, § 1.3{]} Литература 1. Б а р е н б л а т т Г. И., Е н т о в В. М., Р ы ж и к В. М. Движение жидкостей и газов в природных пластах. М.: Недра, 1984. 208 с. (http://gidropraktikum.narod.ru/Barenblatt-et-al-1984.pdf) 2. Б а р е н б л а т т Г. И., Е н т о в В. М., Р ы ж и к В. М. Теория нестационарной фильтрации жидкости и газа. М.: Недра, 1972. 288 с. (http://gidropraktikum.narod.ru/Barenblatt-et-al-1972.djvu) 3. Б а с н и е в К. С., В л а с о в А. М., К о ч и н а И. Н., М а к с и м о в В. М. Подземная гидравлика. М.: Недра, 1986. 303 с. (http://gidropraktikum.narod.ru/Basniev-et-al.djvu) 4. К о л л и н з Р. Течения жидкостей через пористые материалы. Пер. с англ. под ред. Г. И. Баренблатта. М.: Мир, 1964. 352 с. (http://gidropraktikum.narod.ru/Collins.djvu) 5. Л е о н т ь е в Н. Е. Основы теории фильтрации. М.: МАКС Пресс, 2017. 88 с. (http://gidropraktikum.narod.ru/Leontiev-2017.pdf) 6. М а с к е т М. Течение однородных жидкостей в пористой среде. М.--Ижевск: Институт компьютерных исследова- ний, 2006. 640 с. (http://gidropraktikum.narod.ru/Muskat-2004.pdf) 7. П о л у б а р и н о в а - К о ч и н а П. Я. Теория движения грунтовых вод. М.: Наука, 1977. 664 с. (http://eqworld.ipmnet.ru/ru/library/books/Kochina1977ru.djvu) 8. Ч а р н ы й И. А. Подземная гидрогазодинамика. М.: Гостоптехиздат, 1963. 396 с. (http://gidropraktikum.narod.ru/Charny.djvu) 9. B e a r J. Dynamics of fluids in porous media. N.Y.: Dover Publications, 1972. 764 p. (http://gidropraktikum.narod.ru/Bear-Dynamics-of-fluids-in-porous-media-1972.djvu) 10. К у л и к о в с к и й А. Г., С в е ш н и к о в а Е. И. Нелинейные волны в упругих телах. М.: Московский лицей, 1998. 412 с. (http://gidropraktikum.narod.ru/Kulikovskii-Sveshnikova.djvu)

% ===== Раздел 5 =====

\section{И. А. Чарный. Подземная гидрогазодинамика. 1963}\label{section-05}

Файл: \texttt{Charny.djvu}.

\begingroup\small\setlength{\parindent}{0pt}\setlength{\parskip}{1pt}
{[}Страница файла 393{]}\par
ОГЛАВЛЕНИЕ \par
Стр. \par
Предисловие 3 \par
Глава I. Основные понятия теории фильтрации 9 \par
§ 1. Пористая среда. Связь скорости фильтрации с  \par
действительной физической скоростью частиц жидкости. Закон Дарси. \par
Коэффициенты фильтрации и проницаемости 9 \par
§ 2. Одномерное течение. Приток несжимаемой жидкости к стоку \par
и источнику на плоскости и в пространстве 17 \par
§ 3. Одномерное установившееся движение однородной  \par
сжимаемой жидкости в трубке тока переменного сечения 24 \par
§ 4. О распределении скоростей в поперечном сечении  \par
фильтрационного потока 27 \par
Глава II, Дифференциальные уравнения теории фильтрации однородной \par
ЖИДКОСТИ 31 \par
§ 1. Дифференциальные уравнения изотермической фильтрации \par
без учета массовых сил 31 \par
\textbar{} 2. Вывод Н. Е. Жуковского дифференциальных уравнений  \par
изотермического движения жидкости в пористой среде с учетом \par
массовых сил из дифференциальных уравнений Эйлера для \par
идеальной жидкости 33 \par
§ 3. Замечания о системе уравнений для общего случая  \par
неизотермической фильтрации 40 \par
§ 4. Закон сохранения массы фильтрационного потока 41 \par
§ 5. Закон изменения количества движения фильтрационного потока 42 \par
§ G. Закон сохранения энергии для фильтрационного потока . . 48 \par
§ 7. Частные случаи 57 \par
§ 8. Примеры практического использования уравнения энергии \par
для неизотермической фильтрации 61 \par
Глава III. Плоские задачи теории фильтрации о притоке к  \par
совершенным скважинам 66 \par
§ 1. Связь теории функции комплексного переменного с плоской \par
задачей теории фильтрации. Функция тока. Комплексный \par
потенциал 66 \par
§ 2. Приток к точечным стокам и источникам на плоскости. Случай \par
равнодебитных стока и источника. Приток к скважине,  \par
эксцентрично расположенной в круговом пласте 72 \par
§ 3. Вывод некоторых формул для притока к скважинам при \par
помощи конформного отображения 78 \par
§ 4. Течение между конфокальными эллипсами 88 \par
\smallskip
\smallskip
\smallskip
{[}Страница файла 394{]}\par
394 \par
Оглавление \par
Стр. \par
§ 5. Приток к скважинам в пласте овальной формы 89 \par
§ 6. Приток к скважинам в круговом пласте при переменном  \par
давлении на контуре питания 92 \par
Глава IV. Практические методы расчета притока к совершенным  \par
скважинам 100 \par
§ 1. Вводные замечания. Потенциал нескольких точечных стоков \par
в неограниченном пласте 100 \par
§ 2. Приток к совершенной скважине в пласте с прямолинейным \par
контуром питания. Метод отражения 103 \par
§ 3. Приток к группе совершенных скважин в пласте с удаленным \par
контуром питания. Потенциал и скорость фильтрации  \par
результирующего течения 106 \par
§ 4. Приток к бесконечным цепочкам и кольцевым батареям  \par
скважин 110 \par
§ 5. Метод эквивалентных фильтрационных сопротивлений . . . 111 \par
Глава V. Приток к несовершенным скважинам при линейном и нелинейном \par
законах фильтрации 117 \par
§ 1. Виды несовершенства скважин. Приток к одной необсаженной \par
скважине с открытым забоем в центре кругового пласта . . 117 \par
§ 2. Интерференция несовершенных скважин. Приведенный радиус \par
несовершенной скважины 124 \par
§ 3. Расчет фильтрационного сопротивления, обусловленного \par
несовершенством скважины в однородно-анизотропном пласте 129 \par
§ 4. Приток к скважинам с двойным видом несовершенства . . . 139 \par
\textbar{} 5. Случай обсаженной скважины, вскрытой небольшим числом \par
перфораций 142 \par
§ 6. Индикаторные кривые дебит --- депрессия для однородной \par
несжимаемой жидкости и для газа при линейном и  \par
нелинейном законах фильтрации. Интерференция скважин при  \par
двучленном законе фильтрации 145 \par
Глава VI. Безнапорное движение жидкости в пористой среде 157 \par
§ 1. Особенности безнапорного движения. Предпосылки  \par
гидравлической теории безнапорного движения 157 \par
\textbar{} 2. Гидравлическая теория безнапорного движения через  \par
прямоугольную перемычку на горизонтальном непроницаемом  \par
основании 161 \par
§ 3. Гидравлическая теория безнапорного притока к совершенной \par
скважине 162 \par
§ 4. Строгое доказательство формул Дюпюи для безнапорного \par
движения через перемычку и притока к скважине. ..... 163 \par
§ 5. Сведение безнапорного движения к равнодебитному напорному 168 \par
§ 6. Интегральные соотношения между фильтрационным расходом \par
и граничными потенциалами при плоском и осесимметричном \par
движении в пласте переменной мощности 172 \par
§ 7. Дифференциальные уравнения гидравлической теории  \par
нестационарной безнапорной фильтрации. Аналогия с  \par
неустановившейся фильтрацией идеального газа 174 \par
Глава VII. Движение и равновесие границы раздел ;■ длух жидкостей \par
в пористой среде . \textquotesingle. 178 \par
§ 1. Введение 178 \par
§ 2. Уравнения движения отмеченных частиц в потоке однородной \par
жидкости 180 \par
\smallskip
\smallskip
\smallskip
{[}Страница файла 395{]}\par
Оглавление \par
395 \par
Стр. \par
§ 3. Расчет скорости вытеснения одной жидкости другой из трубок \par
тока, предполагаемых неизменными 181 \par
§ 4. Замечания о задачах фильтрации с подвижными граничными \par
условиями. Сведение задачи о движении границы раздела \par
к решению интегро-дифференциального уравнения  \par
специального типа 186 \par
§ 5. Вытеснение нефти водой из трубки тока переменного сечения 193 \par
§ 6. Прямолинейное и плоско-радиальное движение границы  \par
раздела в пласте с постоянными мощностью, пористостью и  \par
проницаемостью 195 \par
§ 7. Характер движения водо-нефтяного контакта. Схемы предельно \par
анизотропных пластов. Устойчивость движения границы \par
раздела 198 \par
§ 8. Конус подошвенной воды. Условия равновесия и прорыва \par
подошвенной воды или верхнего газа в скважину 202 \par
§ 9. Совместный приток нефти и подошвенной воды к  \par
несовершенной скважине 214 \par
■§ 10. Выравнивание возмущенной в начале границы раздела двух \par
жидкостей в пористой среде под действием гравитационных \par
сил. Расчет предельных схем течения 218 \par
§ 11. Автомодельные решения задачи о прямолинейном и  \par
радиальном вытеснении одной жидкости другой в предельно  \par
анизотропном пласте ку = оо 230 \par
Глава VIII. Нестационарная фильтрация однородной упругой жидкости \par
и газа \textquotesingle{} . . . 250 \par
§ 1. Уравнения движения упругой жидкости в упругой пористой \par
среде. Влияние неподвижных газовых включений на  \par
приведенный модуль упругости пластовой системы 250 \par
§ 2. Вывод формул для притока упругой жидкости к  \par
прямолинейной галерее и к точечному источнику на плоскости .... 258 \par
§ 3. Приток к точечному стоку и кольцевой галерее при  \par
переменном дебите 269 \par
§ 4. Метод последовательной смены стационарных состояний.  \par
Приток к прямолинейной галерее 272 \par
§ 5. Расчет радиального притока упругой жидкости по методу \par
последовательной смены стационарных состояний 279 \par
\textbar{} 6. Об одном видоизменении метода интегральных соотношений \par
для решения задач упругого режима фильтрации с  \par
неподвижными а подвижными граничными условиями 284 \par
§ 7. Некоторые методы определения параметров пластов и  \par
скважин при упругом режиме фильтрации 298 \par
\textbar{} 8. Краткие сведения о нестационарной фильтрации газов . . . 306 \par
Глава IX. Фильтрация смесей нескольких жидкостей 313 \par
§ 1. Вводные замечания. Основные уравнения фильтрации  \par
двухфазной жидкости 313 \par
§ 2. Теория Баклея --- Леверетта 321 \par
§ 3. Общие соображения о движении скачков насыщенности . . 327 \par
§ 4. Одномерная фильтрация двухфазной жидкости в трубке тока \par
переменного сечения без учета массовых сил 330 \par
§ 5. Одномерная фильтрация двухфазной жидкости с  \par
произвольными массовыми силами 335 \par
§ 6. Фильтрация двухфазной жидкости с учетом капиллярного \par
давления \textquotesingle. 345 \par
\smallskip
\smallskip
\smallskip
{[}Страница файла 396{]}\par
396 \par
Оглавление \par
Стр. \par
§ 7. Фильтрация трехфазной смеси 348 \par
§ 8. Установившееся движение газированной жидкости в пористой \par
среде. Замечания о практических методах расчета  \par
неустановившегося движения 355 \par
§ 9. Уравнения движения многокомпонентных смесей с учетом \par
фазовых превращений 365 \par
Глава X. Некоторые специальные задачи вытеснения в пористой среде 371 \par
§ 1. Влияние фазовых проницаемостей на характер движения при \par
интенсивном темпе вытеснения одной жидкости другой в  \par
пористой среде 371 \par
§ 2. Влияние неоднородности пористой среды вдоль вертикали \par
на процесс вытеснения 374 \par
§ 3. Расчет движения газового объема в неограниченном  \par
водоносном пласте 378 \par
Л в т о р \par
Исаак Абрамович Чарный \par
ПОДЗЕМНАЯ ГИДРОГАЗОДИНАМИКА \par
Ведущий редактор С. М. каешкова \par
Технический редактор А. С. Полосина. Корректор Л. В. Сметанина \par
Подписано к набору 20/VII 1963 г. Подписано к печати 3/Х 1963 г. \par
Формат бОХЭО\textquotesingle/гв- Печ. л. 24,75- Уч.-изд. л. 26,74. Т-12457. \par
Тираж 2960 экз. Зак. 805/508. Цена 1 р. 09 к. \par
Гостоптехиздат. Москва, К-12, Третьяковский проезд, 1/19- \par
Типография «Красный Печатник». Ленинград, Московский проспект, 91.\par
\endgroup

% ===== Раздел 6 =====

\section{А. Г. Куликовский, Е. И. Свешникова. Нелинейные волны в упругих телах. 1998}\label{section-06}

Файл: \texttt{Kulikovskii-Sveshnikova.djvu}.

\begingroup\small\setlength{\parindent}{0pt}\setlength{\parskip}{1pt}
{[}Страница файла 413{]}\par
Содержание \par
Предисловие 7 \par
Математическое введение 13 \par
1.1. Законы сохранения и соответствующие им \par
дифференциальные уравнения 13 \par
1.2. Гиперболические системы. Линейные и \par
линеаризованные уравнения. Слабые разрывы. \par
Инварианты Римана 18 \par
1.3. Граничные условия. Эволюционность 26 \par
1.4. Волны Римана 31 \par
1.5. Разрывы и соотношения на них. Ударная адиабата 38 \par
1.6. Условия эволюционности разрывов 43 \par
1.7. Разрывы малой интенсивности 50 \par
1.8. Поведение ударной адиабаты в окрестности точек \par
Жуге 55 \par
1.9. Распад произвольного разрыва и другие автомо- \par
дельные задачи 61 \par
1.10. Уравнения законов сохранения в форме Годунова. \par
Энтропия 71 \par
1.11. Решения с разрывами как предел непрерывных \par
решений уравнений усложненной модели 77 \par
1.12. Малые возмущения в среде с диссипацией 81 \par
1.13. Структура ударных волн 84 \par
1.14. Дополнительные соотношения на разрыве 88 \par
1.15. Еще о структуре и о числе дополнительных \par
соотношений .......\textquotesingle{} 96 \par
\smallskip
\smallskip
\smallskip
{[}Страница файла 414{]}\par
4 СОДЕРЖАНИЕ \par
2. Постановка задачи о плоских волнах \par
в упругой среде 117 \par
2.1. Модель упругой среды. Система уравнений 117 \par
2.2. Уравнения теории упругости для одномерных \par
движений в виде плоских волн. \par
Условия на разрыве 126 \par
2.3. Волновая изотропия и анизотропия. Внутренняя \par
энергия среды с малой волновой анизотропией . . .131 \par
2.4. Задание упругого потенциала для \par
слабонелинейной среды 134 \par
2.5. Распространение нелинейных волн в средах, \par
взаимодействующих с электромагнитным полем . . 142 \par
3. Волны Римана 153 \par
3.1. Малые возмущения. Линейные волны 153 \par
3.2. Уравнения для волн Римана 159 \par
3.3. Квазипродольные волны Римана 161 \par
3.4. Квазипоперечные волны Римана 164 \par
3.5. Изменение величин в квазипоперечных волнах . . . 169 \par
3.6. Эволюция квазипоперечных волн Римана 171 \par
3.7. Волны Римана в случае волновой изотропии .... 173 \par
4. Ударные волны 177 \par
4.1. Условия на слабых ударных волнах 177 \par
4.2. Квазипродольные ударные волны 181 \par
4.3. Квазипоперечные ударные волны. Ударная адиабата183 \par
4.4. Условие неубывания энтропии 190 \par
4.5. Условия эволюционности скачка. Точки Жуге . . . 192 \par
4.6. Скорость квазипоперечных ударных волн 194 \par
4.7. Типы эволюционных ударных, волн и их количество 199 \par
4.8. Положение участков эволюционности на ударной \par
адиабате 204 \par
4.9. Ударные переходы в заданное состояние 210 \par
4.10. Частные виды начальной деформации 211 \par
4.11. Квазипоперечные ударные волны при G/R? \textless С 1 \par
и G/R2 \textgreater{} 1 223 \par
4.12. Волны в несжимаемых средах 230 \par
\smallskip
\smallskip
\smallskip
{[}Страница файла 415{]}\par
СОДЕРЖАНИЕ 5 \par
4.13. Представление неэволюционного разрыва двумя эво- \par
люционными, идущими с одинаковой скоростью . . 232 \par
5. Нестационарные автомодельные задачи \par
для волн малой амплитуды 239 \par
5.1. Предварительные замечания 239 \par
5.2. Задача о внезапном изменении нагрузки на \par
границе полупространства ("задача о поршне") . .241 \par
5.3. Построение решения в средах сх\textgreater0 245 \par
5.4. Построение решения в средах сх\textless0 265 \par
5.5. Распад произвольного начального разрыва. \par
Соударение ударных волн 273 \par
6. Двумерные стационарные нелинейные волны \par
и автомодельные решения 283 \par
6.1. Простые волны 283 \par
6.2. Стационарные ударные волны в движущейся среде 290 \par
6.3. Задача о бегущей нагрузке 294 \par
7. Упрощенные уравнения теории упругости. \par
Неавтомодельные задачи 297 \par
7.1. Эквивалентная несжимаемая среда при описании \par
квазипоперечных волн 297 \par
7.2. Упрощенные уравнения для квазипоперечных волн, \par
распространяющихся в одну сторону 301 \par
7.3. Описание волн Римана и разрывов с помощью упро- \par
щенных уравнений теории упругости 304 \par
7.4. Инварианты Римана. Уравнения в плоскости \par
годографа. Неавтомодельные задачи 306 \par
7.5. Эволюция при малой анизотропии волн, близких к \par
вращательным. Катастрофическая перестройка . . 310 \par
8. Движения вязкоупругих сред 317 \par
8.1. Модель Кельвина-Фойхта. Уравнения одномерных \par
движений 317 \par
8.2. Задача о структуре квазипоперечных ударных волн 320 \par
\smallskip
\smallskip
\smallskip
{[}Страница файла 416{]}\par
6 СОДЕРЖАНИЕ \par
8.3. Численные решения задач о вязкоупругих волнах, \par
имеющие автомодельную асимптотику 333 \par
9. Волны конечной амплитуды в слабоанизотропных \par
упругих средах 359 \par
9.1. Введение. Случай волновой изотропии 359 \par
9.2. Волны Римана в несжимаемой среде при д \textless С 1 . . . 365 \par
9.3. Волны Римана в несжимаемой среде при \par
анизотропии частного вида 374 \par
9.4. Волны Римана в сжимаемой среде 377 \par
9.5. Ударные волны в несжимаемой среде 381 \par
9.6. Ударные волны в сжимаемой среде 392 \par
Литература 399 \par
Предметный указатель 407\par
\endgroup

% ===== Раздел 7 =====

\section{R. E. Collins. Flow of Fluids Through Porous Materials. Английское издание}\label{section-07}

Файл: \texttt{Collins-eng.pdf}.

\begingroup\small\setlength{\parindent}{0pt}\setlength{\parskip}{1pt}
{[}Страница файла 7{]}\par
6 PETROMAN\par
\smallskip
CONTENTS\par
\smallskip
Preface. .\par
\smallskip
Chapter 1. Structure and Properties of Porous Materials 1\par
1.10 Structure and Classification. .......0....6 06sec 1\par
1.20 Structure and Properties 2\par
1.30 Porosity coe 3\par
1,31 Porosity and Structure... 3\par
1.32 Fffects of Consolidation and Compaction on Porosity. 5\par
1.33; Measurement of Porosity cece 7\par
\smallskip
1.40 Specific Surface\par
\smallskip
1.41 Measurement of Specific Surface... =... : 9\par
1.50 Permeability : see 10\par
1,51 Structural Interpretation of Permeability. . as .. it\par
\smallskip
1.52 Factors Affecting the Permeability of Porous Materials B\par
1.60 `Typical Values of Porosity, Pormeability and Specific Surfiee. 14\par
1,70 Porosity and Permeability Distributions in Naturally Porous\par
\smallskip
Materiais : . 4\par
1.80 Mechanical Properties of Porous Materials....00..0..0... 0 19\par
Exercises... re 20\par
Referene a\par
Chapter 2. Statics of Fluids in Porous Media 22\par
2.10 Fluid Saturations. . . eee cece eee 22\par
2.11 Measurement of Fluid Saturations.. 0.0000. 0cs.c sees BB\par
2.20 Capillary Pressure... beens 23\par
2.21 Capillary Pressure Saturation Relationship... 25\par
2.22 Measurement of Capillary Pressure........000.00.600000.. 26\par
2.23 Capillary Hysteresis............. 29\par
\smallskip
2.24 Lrreducible Saturation (Connate Water)\par
2.25 The Leverett j-Function......... .\par
2.30 Pore-Size Distribution..........2.... .\par
2.40 Vapor Pressuro of Fluids in Porous Media...\par
\smallskip
2.41 Measurement of Vapor Pressure-Saturation Curves......... 35\par
2.50 Heat of Wetting and Heat of Swelling. .... ve BG\par
2.60 Electrical Properties of Fluid-filed Porous Materials........ 38\par
2.70 Osmosis... 0.06.6 ee renee ceeeeeese ees 41\par
\smallskip
vii\par
\smallskip
www.petroman.ir\par
\smallskip
\smallskip
{[}Страница файла 8{]}\par
2.80 Electrochemical Potentials. 00.0... AD\par
\{ 2.0 Static Muid Distributions in Porous Materials, . 4\par
i Exel 6\par
i References 46\par
\textbar{} 3. Physical and Mathematical Theory of Flow 47\par
\{ Mechanisms and Types of Flow............. 47\par
j Laminar Viscous Flow through Porous Media.............. 47\par
1 `Turbulent Plow through Porous Media cette BL\par
` Simultaneous Laminar Mow of Immiscible Fluids through\par
j Porous Media. 0.0.0 ecco : fevers BB\par
Streaming Potential and Electro-Osmosis.....-.- 56\par
5 Thermo-Osmosis, renner 58\par
Limitations of Classical Hydrodynamics... 00.58\par
Differential Form of Darcy\textquotesingle s Law. Leteeeee B\par
Simultaneous Flow of Immiscible Fluids through Porous\par
Media; Relative Permeability \_ 60\par
Darcy\textquotesingle s Law for Anisotropic Porous Media 2\par
Equations of State for Muids. .. veteeeeeees 65\par
`The Continuity Equation. . fie eeeee 66\par
The Differential Equations of Fluid Flow through Porous\par
Materials. i vines 69\par
3.70 Boundary Conditions. 74\par
Referene: feces ces cece 16\par
Chapter 4. Steady Laminar Flow of Homogeneous Fluids a7\par
4.10 The Steady-State Regime Li ceeceeeeeeeeeee WT\par
4.20 Linear Plow and the Measurement of Permeability... 77\par
4.30 Horizontal Plane Flow in Homogeneous Media; Mathematical\par
Vormutation........ Le tteeeeeeee. 80\par
`The Stream Function and Complex Potential ceeeee 82\par
Radial Flow between Coneentric Circular Cylinders... 86\par
: Conformal Mapping. fo ieteecteteeseeees) 88\par
`The Schwarz-Christoffel Transformation; Conformal Mapping\par
of a Rectangle. . ee ceeceeees 91\par
440 `Transverse Permeabilities of Cylindrical Cores... ., 93\par
4.50 Sourees and Sinks een Letcceereeees 98\par
4.51 System of Wells Near a Plane Discontinuity in Medium; the\par
Method of Images. .... coves fovveeeeecees 102\par
4.60 Gravity Drainage: Free Surfaces... 105\par
\smallskip
www.petroman.ir\par
\smallskip
{[}Страница файла 9{]}\par
6 PETROMAN conrenrs ix\par
\smallskip
Bxereises... 0... ce 106\par
References. . . . =e ene .. 107\par
\smallskip
Chapter 5. Transient Laminar Flow of Homogeneous Fluids\par
5.10 Transient Flow of Compressible Liquids\par
5.20 The Laplace Transform and Duhamel's Theorem\par
5.30 Sources and Sinks; the Method of Images.\par
5.40 Pressure Build-up Tests in Oil Wells.\par
5.50 Performance of Water-Drive Reservoirs\par
51 The Infinite Linear Aquifer.....\par
52 The Infinte Radial Aquifer\par
5,53 The Tilted Aquifer. wee\par
5.60 Flow of an Ideal Gas, Numerical Integration\par
\smallskip
Chapter 6, Simultaneous Flow of immiscible Fluids 139\par
6.10 Steady Linear Flow and the Measurement of Relative Perme-\par
ability; the Boundary Hilfect\par
6.20 Immiscible Displacement; The Buckley-Leverett Equation\par
6.21 The Welge Integration of the Bucktey-Leverett Equation; The\par
Calculation of Relative Permeabilities from Displacement\par
\smallskip
Data...... 149\par
6.30 Linear Immiscible Displacement Including the Effects of\par
Capillary Pressure eee bes 182\par
640 Imbibition Si aesseterenesee . 161\par
50 Multi-dimensional `Two-phase Flow: Numerical Solutions... 167\par
6.60 Turbulent Flow of Immiscible Fluids. . . 168\par
Exercises. . ae \_ eee a 168\par
References rn eee eee oe ... 169\par
Chapter 7. Moving Boundary Problems; Displacement; Deposition of\par
Solids 170\par
7.10 Displacement of One Fluid by Another 170\par
7.11 Analytical Solution of the Linear Displacement. Problem 73\par
\smallskip
7.20 Two-Dimensional Displacement and Waterflood Sweep Pat\par
terns for Unit Mobility Ratio 115\par
7.21 The Effect of Mobility-Ratio on Areal Sweep Efficiency... 187\par
\smallskip
7.30 Filtration and the Deposition of Solids........ veces 189\par
740 Frontal Instability and Viscous Fingering........0........ 196\par
Exercises. . oe . . 200\par
References . 200\par
\smallskip
www.petroman.ir\par
\smallskip
\smallskip
{[}Страница файла 10{]}\par
6 PETROMAN\par
\smallskip
x CONTENTS\par
\smallskip
Chapter 8.\par
\smallskip
imultaneous Laminar Flow of Miscible Fluids 201\par
\smallskip
8.10 Miscible and Immiscible Fluids, Fick's Law of Diffusion 201\par
8.20 Miscible Displacement in a Capillary Tube........ ss. 202\par
8.30 Miscible Displacement in Porous Media........ wee 207\par
8.31 Theory of Dispersion in Porous Media................. 200\par
\smallskip
8.40 Macroscopic Features of Miscible Displacement.\par
8.50 Turbulent Flow of Miscible Fluids.\par
\smallskip
oe. 212\par
.. 214\par
. 25\par
\smallskip
Exercises... OD e eee cece eee eee ete eee e ees .\par
References. 0.2.2. e cece eee cece eres QU5\par
Chapter 9. Theory of Models 217\par
9.10 The feeb e eee e eee se. 217\par
9.20 Basic Theory of Models eee vee 218,\par
9.30 Scaled Models and Scaling Laws. teens » 221\par
\smallskip
9.40 Examples of Scaling Analy:\par
\smallskip
System are Known. ... wees weeee 222\par
9.50 Example of Analog Construction... 229\par
9.60 Example of Scaling When the Bauations for the System are\par
\smallskip
not Known... 2.0.2... eeee ee . ve. 282\par
9.70 Partial Similarity in Sealed Models...... veceeeee BBB\par
\smallskip
9.80 Special Aspects of Model Scaling for Porous Media . . 238\par
9.90 Results of Model Studies of Flow Through Porous Media... 239\par
\smallskip
Appendix... . 240\par
Exercises. ce . . cece 240\par
References. .......... vevee eee eeeeeees cee SAL\par
Chapter 10. Flow with Change of Phase 242\par
10.10 Flow With Change of Phase; Solution-Gas Drive........... 242\par
10.20 Phase Equilibrium. . bet eecer eee 2. 242\par
10.30 Supersaturation and Undersaturation........ ss... 50\par
10.40 Formulation of the Solution-Gas Drive Process...... 253\par
10.50 Second Formulation of the Solution-Gas Drive Process...... 256\par
10.60 Third Formulation of the Solution-Gas Drive Process; The\par
Material Balance Equations. ... + 259\par
10.70 Flow with Change of Phase; General Procedures. . » 262\par
Bxxereises..0. eee cee ee + 263\par
References... + 263\par
Index....,. » 264\par
\smallskip
www.petroman.ir\par
\endgroup

% ===== Раздел 8 =====

\section{Г. И. Баренблатт, В. М. Ентов, В. М. Рыжик. Теория нестационарной фильтрации жидкости и газа. 1972}\label{section-08}

Файл: \texttt{Barenblatt-et-al-1972.djvu}.

\begingroup\small\setlength{\parindent}{0pt}\setlength{\parskip}{1pt}
{[}Страница файла 289{]}\par
ОГЛАВЛЕНИЕ\par
\smallskip
Глава Г. Физические основы .... ще" + sae\par
\$1. Пористая среда и ee свойства зе -24.- К -\par
\$ 2. Закон фильтрации однородной жидкости \textless{} .\par
\$ 3. Зависимость параметров жидкости и порвотой среды от\par
давления еее eee\par
\smallskip
Глава И. Осповные задачи нестационарной фильтрации... .....\par
\smallskip
\$1. Уравнение неразрывности „........-\par
\smallskip
§ 2. Упругий режим фильтрации... еее\par
\$ 3. Уравнения безпапорно "фильтрации несжимаемой жидкости\par
\$4. Основные уравнения фильтрации газа... .....\par
\smallskip
Глава ИТ. Теория упругого режима фильтрации... ..-\par
\smallskip
§ 1. Одномерное прямолипейно-параллельное движение ....\par
\$2. Осесимметричные задачи и задачи интерференции скважин\par
при неустановившейся фильтрации . .\par
\smallskip
\$ 3. Некоторые специальные задачи упругого режнма ....\par
\$ 4. Обратные задачи линейной теории нестационарной фильтрации\par
\smallskip
Глава ГУ. Нелинейные инвариантные задачи нестационарной фильтрации\par
жидкостей и газов „еее ..\par
\smallskip
\$1. Общая характеристика инвариантных задач теории нестацио-\par
нарной фильтрации. Автомодельные пологие безнапорные дви-\par
жения при нулевом начальном уровне жидкости .\par
\$ 2. Пологие безнапорные движения с нулевым начальным напором:\par
предельные автомодельные движения, оевсииметрачные авто-\par
модельные движения 6 6 ee ee ee .\par
\$3. Артомодельвые движения жидкости H газа плоскими волнами\par
в полубесконечиом пласте при ненулевом начальном давлении\par
таза или уровне жидкости... .. иене\par
\$ 4. Осесимметричные автомодельные движения в бесконечном пласте\par
при непулевом начальном давлении газа или уровне жидкости\par
\$ 5. Некоторые специальные автомодельные движения ...\par
\smallskip
Глава У. Приближенпыс методы решения задач нестационарной Фильтра-\par
\smallskip
\$1. Схема метода интегральных соотношений. Приток к талорее\par
в бесконечном пласте при упругом режиме. .\par
\smallskip
§ 2. Решение задач упругого режима методом `интегральных соотно-\par
шений .. toe coe tee\par
\smallskip
§ 3. Решение задач HecraumonapHoll фильтрации raga...\par
\smallskip
23\par
\smallskip
26\par
32\par
37\par
47\par
\smallskip
116\par
125\par
\smallskip
{[}Страница файла 290{]}\par
\$ 4. Видоизменоние метода интегральных соотношений для случая\par
бесконечного пласта еее еее еее 10\par
\smallskip
Глава УТ. Фильтрация неоднородных жидкостей... :... 146\par
\smallskip
\$1. Фильтрация многофазных жидкостей Чена 146\par
§ 2. Вытеснение несмешивающихся жидкостей. Задача Баклея ---\par
\smallskip
«Леверетта еее еее. 155\par
\$ 3. Структура скачка насыщенности. Уравнение Рапопорта --- Лиса.\par
\smallskip
Стабилизированная зона. еее -е. 464\par
\$ 4. Капиллярная пропитка и автомодельные задачи вытеснения ве-\par
\smallskip
смешивающихся жидкостей „еее неее 112\par
\smallskip
Глава УИ. Неуставовившаяся фильтрация в  трощиновато-пористых\par
средах, средах с двойной пориетостыо и слоистых породах 185\par
\smallskip
\$ 1. Фильтрация однородной жидкости в трещиновато-порнстой среде 185\par
\$ 2. Основные задачи неустановившейся фильтрации однородной\par
\smallskip
жидкости в трещиноватых и слоистых пластах... ..... 19\%\par
\smallskip
\$3. Двухфазная нестационарная фильтрация и вытеснение не-\par
смешивающихся жидкостей в средах с двойной пористостью .. 202\par
Глава УИТ. Нелинейная исстационарная фильтрация... ...... 217\par
\smallskip
§ 1. Отклонения ог закона Дарси. Нелинейная фильтрация .... 27\par
\$2. Автомодельные решения залачи фильтрации жидкости при не-\par
`линейном законе сопротивления. Фильтрация с предельным гра-\par
\smallskip
\$ 3. Приближениое решение задач нестационарной фильтрации при\par
велинейпом законе сопротивления ен. 28\par
\smallskip
Глава IX, Специальные задачи нестационарной фильтрации однородной\par
\$ 1. Неизотермическое фильтрационное дви ю ee ee BBB\par
\$2. Упруго-иластический режим нефтяного miacta. . . 0... 25\par
\smallskip
Глава Х. Специальные задачи нестационарной фильтрации неоднородной\par
жидкости. еее иене . 255\par
§ 1. Вытеснение взаимосмешивающихся жидкостей... .... 255\par
\smallskip
\$2. Устойчивость вытеснения несмешивающихся и смешивающихся\par
жидкостей из пористой среды se ee ee ee 261\par
Приложение еее еее еее 2\par
Список литературы еее еее еее еее се 28\par
\smallskip
Григорий Исаакович Баренблатт, Владимир Мордутовим Ентов,\par
Виктор Мизайлович Рыжик\par
\smallskip
ТЕОРИЯ НЕСТАЦИОНАРНОЙ ФИЛЬТРАЦИИ ЖИДКОСТИ И ГАЗА.\par
\smallskip
Редактор Н. Д. Дубровина Технический редактор В. В. Соколова\par
рН. д. Huse Корректор М. П. Курыяма\par
Сдано в набор 13/УИ 1972 г. Подписано в печать 17/Х 1972г. Т-17124.\par
Формат 60X90 ie. Бумага № 2. Печ. л. 18.0. `Уч.-изд. л. 18,54.\par
\smallskip
Тираж 1500 экз. Заказ № 1865/72-6. Цена 2 р. 08 к.\par
\smallskip
`Издительство «Недра», 103633, Москва, К-12, Третьяковокий проезд, 1/19.\par
`Пепингралокая типография NM 6 Главполиграфирома Государственного\par
Комитета Совета Министров СССР по делам издательств, полиграфии и книжной торговли.\par
"Московский пр., 94.\par
\endgroup

% ===== Раздел 9 =====

\section{К. С. Басниев, А. М. Власов, И. Н. Кочина, В. М. Максимов. Подземная гидравлика. 1986}\label{section-09}

Файл: \texttt{Basniev-et-al.djvu}.

\begingroup\small\setlength{\parindent}{0pt}\setlength{\parskip}{1pt}
{[}Страница файла 302{]}\par
ОГЛАВЛЕНИЕ \par
Введение 3 \par
Глава 1. Основные понятия и законы фильтрации нефти, газа и \par
воды 6 \par
§ 1. Особенности движения флюидов в природных пластах 6 \par
§ 2. Пористая среда и ее фильтрационные характеристики 6 \par
§ 3. Скорость фильтрации. Закон Дарси --- линейный  \par
закон фильтрации 9 \par
§ 4. Границы применимости закона Дарси. Нелинейные \par
законы фильтрации 13 \par
§ 5. Обобщенный закон Дарси 19 \par
§ 6. Понятие о режимах нефтегазоводоносных пластов . 22 \par
Глава 2. Основы моделирования процессов фильтрации нефти, газа \par
и воды \textquotesingle. 24 \par
§ 1. Цели и задачи моделирования фильтрационных  \par
процессов 24 \par
§ 2. Физическое моделирование процессов фильтрации  \par
пластовых флюидов 25 \par
§ 3. Математическое моделирование процессов фильтрации \par
пластовых флюидов 26 \par
§ 4. Численная и машинная модели пластовых систем ... 27 \par
§ 5. Аналоговое моделирование 27 \par
§ 6. Основы анализа размерностей и теории подобия ... 28 \par
§ 7. Применение методов теории размерностей в подземной \par
гидравлике 36 \par
Глава 3. Дифференциальные уравнения фильтрации флюидов в  \par
нефтегазоносных пластах 39 \par
§ 1. Вводные замечания 39 \par
§ 2. Вывод уравнения неразрывности 40 \par
§ 3. Дифференциальные уравнения движения 42 \par
§ 4. Зависимость параметров флюидов и пористой среды от \par
давления 45 \par
§ 5. Начальные и граничные условия 49 \par
Глава 4. Установившееся движение несжимаемой жидкости в  \par
пористой среде 51 \par
§ 1. Вывод дифференциального уравнения установившейся \par
фильтрации несжимаемой жидкости по закону Дарси .... 51 \par
§ 2. Одномерные фильтрационные потоки несжимаемой  \par
жидкости в однородных пластах 53 \par
§ 3. Исследование фильтрационных течений несжимаемой \par
жидкости в неоднородных пластах 69 \par
§ 4. Одномерные фильтрационные потоки несжимаемой  \par
жидкости при нелинейных законах фильтрации 78 \par
§ 5. Интерференция скважин . 82 \par
§ 6. Приток к несовершенным скважинам 96 \par
§ 7. Установившиеся безнапорные фильтрационные потоки \par
жидкости 101 \par
§ 8. Решение плоских задач фильтрации методами теории \par
функций комплексного переменного . . • 105 \par
301 \par
\smallskip
\smallskip
\smallskip
{[}Страница файла 303{]}\par
Глава 5. Установившееся движение упругой жидкости и газа в пористой \par
среде . . . ." 11S \par
§ 1. Дифференциальные уравнения установившейся  \par
фильтрации упругой жидкости и газа по закону Дарси ИЗ \par
§ 2. Аналогия установившейся фильтрации сжимаемого  \par
флюида с фильтрацией несжимаемой жидкости 115 \par
§ 3. Установившаяся фильтрация упругой жидкости . . . .116 \par
§ 4. Прямолинейно-параллельный фильтрационный поток \par
идеального газа 117" \par
§ 5. Плоскорадиальный фильтрационный поток идеального \par
газа по закону Дарси 119 \par
§ 6. Плоскорадиальный фильтрационный поток идеального \par
газа по двучленному закону фильтрации 123- \par
§ 7. Плоскорадиальный фильтрационный поток реального \par
газа по закону Дарси 125 \par
§ 8. Фильтрационный поток реального газа по двучленному \par
закону фильтрации к несовершенной скважине 126 \par
Глава 6. Неустановившееся движение упругой жидкости в пористой \par
среде 12S \par
§ 1. Упругий режим пласта и его характерные особенности . 12\& \par
§ 2. Подсчет упругого запаса жидкости в пласте 130 \par
§ 3. Дифференциальное уравнение неустановившейся  \par
фильтрации упругой жидкости 131 \par
§ 4. Одномерные фильтрационные потоки упругой жидкости. \par
Точные решения уравнения пьезопроводности. Основная  \par
формула теории упругого режима 133- \par
§ 5. Интерференция скважин в условиях упругого режима . 144 \par
§ 6. Определение коллекторских свойств пласта по данным \par
исследования скважин при упругом режиме 148- \par
§ 7. Приближенные методы решения задач теории упругого \par
режима 151 \par
§ 8. Приток упругой жидкости к укрупненной скважине . . 163- \par
Глава 7. Неустановившееся движение газа в пористой среде . . . .170 \par
§ 1. Вывод дифференциального уравнения Лейбензона . . . 170 \par
§ 2. Линеаризация уравнения Лейбензона и основное решение \par
линеаризованного уравнения 172 \par
§ 3. Решение автомодельной задачи о притоке газа к скважине \par
с постоянным дебитом в точной постановке 175 \par
§ 4. Решение задачи о притоке газа к скважине методом  \par
последовательной смены стационарных состояний 177 \par
§ 5. Метод усреднения 180 \par
§ 6. Применение принципа суперпозиции к задачам  \par
неустановившейся фильтрации газа 182 \par
§ 7. Приближенное решение задач об отборе газа из  \par
замкнутого пласта при помощи уравнения материального баланса . 185 \par
Глава 8. Взаимное вытеснение жидкостей и газов. Задачи с  \par
подвижными границами 187 \par
§ 1. Кинематические условия на подвижной границе раздела \par
при взаимном вытеснении жидкостей 188 \par
§ 2. Прямолинейно-параллельное вытеснение нефти водой . . 180 \par
§ 3. Плоскорадиальное вытеснение нефти водой 193 \par
§ 4. Устойчивость движения границы раздела жидкостей . . 196- \par
Глава 9. Основы теории фильтрации многофазных систем 198 \par
§ 1. Связь с проблемой нефтегазоотдачи пластов 198 \par
302 \par
\smallskip
\smallskip
\smallskip
{[}Страница файла 304{]}\par
§ 2. Основные характеристики многофазной фильтрации . . 199 \par
§ 3. Исходные уравнения многофазной фильтрации .... 202 \par
§ 4. Одномерные модели вытеснения несмешивающихся  \par
жидкостей 205 \par
§ 5. Задача Баклея---Леверетта и ее обобщения 210 \par
§ 6. Двухфазное течение с учетом капиллярного давления . 222 \par
§ 7. Фильтрация трехфазной смеси 225 \par
§ 8. Движение газированной жидкости в пористой среде . 228 \par
Глава 10. Гидродинамические модели методов повышения нефте- и \par
згазоконденсатоотдачи пластов 232 \par
§ 1. Физические основы применения химических реагентов \par
при заводнении нефтяных пластов 232 \par
§ 2. Вывод системы уравнений двухфазной фильтрации с  \par
активной примесью 234 \par
§ 3. Движение скачков насыщенности и концентрации . . . 236 \par
§ 4. Вытеснение нефти раствором активной примеси .... 238 .. \par
§ 5. Основы неизотермической фильтрации жидкостей и газов 244 \par
§ 6. Температурное поле пласта. Определение потерь теплоты \par
через кровлю и подошву 246 \par
§ 7. Нагревание призабойной зоны пласта при закачке  \par
горячей жидкости 249 \par
Глава 11. Особенности фильтрации неньютоновских жидкостей . . 251 \par
§ 1. Реологические модели фильтрующихся жидкостей и  \par
нелинейные законы фильтрации 252 \par
§ 2. Одномерные задачи фильтрации вязкопластичной  \par
жидкости 257 \par
§ 3. Образование застойных зон при вытеснении нефти водой 262 \par
Глава 12. Движение жидкостей и газов в трещиноватых и трещинс- \par
вато-пористых средах 263 \par
§ 1. Особенности фильтрации в трещиноватых и трещиновато- \par
пористых средах : 263 \par
§ 2. Вывод дифференциальных уравнений движения жидкости \par
и газа в трещиноватых и трещиновато-пористых средах . . . 268 \par
§ 3. Установившаяся одномерная фильтрация жидкости и газа \par
в трещиноватом и трещиновато-пористом пласте 269 \par
§ 4. Неустановившееся движение жидкости и газа в  \par
трещиноватых и трещиновато-пористых средах 274 \par
§ 5. Вытеснение нефти водой из трещиновато-пористых и  \par
неоднородных сред 276 \par
Глава 13. Численное моделирование основных процессов фильтрации \par
пластовых флюидов 279 \par
§ 1. Методы решения задач подземной гидравлики 279 \par
§ 2. Построение дискретного аналога области фильтрации . . 280 \par
§ 3. Разностные схемы, аппроксимирующие уравнения  \par
фильтрации упругой жидкости 281 \par
§ 4. Решение системы конечно-разностных уравнений.  \par
Сходимость метода 285 \par
§ 5. Численное решение двумерной задачи нестационарной \par
фильтрации упругой жидкости в неоднородном пласте .... 288 \par
§ 6. Построение разностных схем для уравнений двухфазной \par
фильтрации несжимаемых жидкостей 292 \par
Список литературы 295 \par
Приложения 296 \par
303 \par
\smallskip
\smallskip
\smallskip
{[}Страница файла 305{]}\par
УЧЕБНИК \par
Каплан Сафербиевич Басниев \par
Анатолий Михайлович Власов \par
Ираида Николаевна Кочина \par
Вячеслав Михайлович Максимов \par
ПОДЗЕМНАЯ ГИДРАВЛИКА \par
Редактор издательства Т. К. Лазарева \par
Технические редакторы О. А. Колотвина, В. Ю. Любимова \par
Корректор К. С. Торопцева \par
ИБ № 5801 \par
Сдано в набор 08.01.86 г. Подписано в печать 29.05.86. Т-07997. Формат \par
60Х90\textquotesingle Дб. Бумага кн.-журн. имп. Гарнитура Литературная. Печать  \par
высокая. Усл. печ. л .19,0. Усл. кр.-отт. 19,0. Уч.-изд. л. 19,0. Тираж \par
4600 экз. Заказ 218/439---6. Цена 95 коп. \par
Ордена «Знак Почета» издательство «Недра», 103633, Москва, \par
Третьяковский проезд, 1/19 \par
Ленинградская типография № 4 ордена Трудового Красного Знамени  \par
Ленинградского объединения «Техническая книга» им. Евгении Соколовой \par
Союзполиграфпрома при Государственном комитете СССР по делам  \par
издательств, полиграфии и книжной торговли. 191126, Ленинград,  \par
Социалистическая, 14\par
\endgroup

% ===== Раздел 10 =====

\section{Н. Е. Леонтьев. Основы теории фильтрации. 2-е изд., 2017}\label{section-10}

Файл: \texttt{Leontiev-2017.pdf}.

\begingroup\small\setlength{\parindent}{0pt}\setlength{\parskip}{1pt}
{[}Страница файла 3{]}\par
Оглавление\par
Предисловие . . . . . . . . . . . . . . . . . . . . . . . . . . . . . . . . .\par
6\par
Глава I.\par
Основные понятия и уравнения теории\par
фильтрации . . . . . . . . . . . . . . . . . . . . . . . . .\par
7\par
§ 1. Описание фильтрации в рамках механики сплош-\par
ных сред . . . . . . . . . . . . . . . . . . . . . . . . . . . . . . .\par
7\par
Фильтрация (7).\par
Описание фильтрации в рамках\par
МСС (7).\par
§ 2. Параметры, описывающие накопление и перенос\par
жидкости в пористой среде . . . . . . . . . . . . . . . . .\par
9\par
Пористость (9).\par
О введении характерного размера\par
пор (10). Просветность (10). Об осредненных геомет-\par
рических параметрах (11). Скорость фильтрации (11).\par
О моделировании структуры пористой среды (13).\par
§ 3. Законы сохранения массы и импульса при филь-\par
трации в пористой среде . . . . . . . . . . . . . . . . . . .\par
13\par
Закон сохранения массы (14).\par
Уравнение баланса им-\par
пульса --- закон Дарси (14). Коэффициент проницаемо-\par
сти (16). Коэффициент фильтрации (16). Закон Дарси\par
для анизотропных сред (17).\par
§ 4. Законы фильтрации, отличные от закона Дарси .\par
18\par
Пределы применимости закона Дарси (18). Вид закона\par
фильтрации при немалых скоростях фильтрации (18).\par
Двучленный закон фильтрации (21). О «турбулентном»\par
режиме фильтрации (21). Фильтрация неньютоновских\par
жидкостей (21).\par
§ 5. Вывод дифференциальных уравнений фильтрации\par
с помощью осреднения микроскопических уравне-\par
ний . . . . . . . . . . . . . . . . . . . . . . . . . . . . . . . . . . .\par
24\par
Осреднение параметров жидкости в пористой среде (24).\par
Свойства операции осреднения (25).\par
Осреднение мик-\par
роскопического уравнения неразрывности (26). Осред-\par
нение микроскопического уравнения движения (26).\par
Упрощение уравнения движения при малых скоро-\par
стях (27).\par
Осреднение микроскопического уравнения\par
энергии (29). Упрощение уравнения энергии (30).\par
§ 6. Соотношения на разрывах и типичные граничные\par
условия . . . . . . . . . . . . . . . . . . . . . . . . . . . . . . . .\par
32\par
Условия на разрывах (32). Непроницаемая граница (33).\par
Граница пористой среды с чистым флюидом (34). Гра-\par
ница раздела фаз внутри пористой среды (34).\par
Про-\par
\smallskip
{[}Страница файла 4{]}\par
Оглавление\par
межуток высачивания (35).\par
Граница двух пористых\par
сред (36).\par
Глава II. Фильтрация несжимаемой жидкости в неде-\par
формируемой пористой среде . . . . . . . . . .\par
38\par
§ 7. Замкнутая система уравнений и простейшие задачи\par
38\par
Замкнутая система уравнений (38).\par
Опыт Дарси (39).\par
Точечный источник в пористой среде (40).\par
Скважина,\par
вскрывающая тонкий пласт (42).\par
§ 8. Применение методов теории функций комплексно-\par
го переменного для решения плоских фильтраци-\par
онных задач . . . . . . . . . . . . . . . . . . . . . . . . . . . .\par
43\par
Плоскопараллельная фильтрация в вертикальной плос-\par
кости (43). Непроницаемая граница (45). Граница с во-\par
доемом (45). Поверхность депрессии (46). Поверхность\par
высачивания (47).\par
Непроницаемая плотина с плоским\par
основанием на пористом слое бесконечной глубины (47).\par
Приток к дренажной щели (50).\par
Течения в плоскости\par
пласта (51).\par
Скважина в поступательном потоке (52).\par
О сходстве и различии фильтрационных течений и по-\par
тенциальных течений идеальной жидкости (53).\par
§ 9. Гидравлическая теория безнапорной фильтрации\par
54\par
Предположения гидравлической теории (54). Уравнение\par
Буссинеска (55). Формулы Дюпюи для безнапорных те-\par
чений (56).\par
Строгое доказательство формулы Дюпюи\par
для фильтрации через прямоугольную плотину (57).\par
§ 10. Аналоговое моделирование фильтрации несжима-\par
емой жидкости . . . . . . . . . . . . . . . . . . . . . . . . . .\par
58\par
Течение в лотке Хеле-Шоу (59).\par
Метод электрогидро-\par
динамической аналогии (62).\par
Глава III. Фильтрационные течения с учетом сжи-\par
маемости флюида и пористой среды . . . .\par
66\par
§ 11. Изотермическая фильтрация газа в недеформиру-\par
емой пористой среде . . . . . . . . . . . . . . . . . . . . . .\par
66\par
Уравнение Лейбензона (66).\par
Граничные условия (67).\par
Автомодельная задача о закачке газа в пласт (68).\par
§ 12. Фильтрация с учетом слабой сжимаемости жидко-\par
сти и пористого скелета . . . . . . . . . . . . . . . . . . . .\par
70\par
Упругий режим фильтрации (70). Распространение вол-\par
ны давления в пласте (72). Пуск скважины с постоян-\par
ным дебитом (73). Об экспериментальном определении\par
4\par
\smallskip
{[}Страница файла 5{]}\par
Оглавление\par
параметров пласта (75). О связи уравнений Лейбензона,\par
пьезопроводности и Лапласа (76).\par
§ 13. Неизотермическая фильтрация газа в недеформи-\par
руемой пористой среде . . . . . . . . . . . . . . . . . . . .\par
77\par
Термодинамика двухпараметрических сред (77).\par
Ана-\par
лиз уравнения баланса энергии (81). Процесс Джоуля ---\par
Томсона (83).\par
Литература . . . . . . . . . . . . . . . . . . . . . . . . . . . . . . . . . .\par
85\par
5\par
\endgroup

% ===== Раздел 11 =====

\section{М. Маскет. Течение однородных жидкостей в пористой среде. 2004}\label{section-11}

Файл: \texttt{Muskat-2004.djvu}.

\begingroup\small\setlength{\parindent}{0pt}\setlength{\parskip}{1pt}
{[}Страница файла 624{]}\par
ОГЛАВЛЕНИЕ \par
Стр, \par
Предисловие 3 \par
Краткая история развития подземной гидравлики в СССР 6 \par
часть первая \par
ОСНОВЫ \par
Глава I. Введение \par
1. Предмет изучения 15 \par
2. Однородные жидкости 16 \par
3. Пористая среда 18 \par
4. Общие положения аналитической теории 19 \par
5. Природа пористой среды . „ 21 \par
6. Систематическая укладка шаров 21 \par
7. Укладка естественных материалов 24 \par
8. Сжимаемость песка и гравия 24 \par
9. Сжимаемость глин 25 \par
10. Влияние уплотнения и деформации отложений на проницаемость . . 26 \par
11. Классификация пород и их свойства по накоплению жидкостей ... 26 \par
12. Связь структуры горных пород и подземных жидкостей . 31 \par
13. Подземные жидкости, их залегание и миграция 34 \par
14. Залегание грунтовых вод 36 \par
15. Содержание воды в несцементированных поверхностных отложениях 37 \par
16. Движение жидкости в капиллярной зоне 37 \par
17. Движение жидкости ниже водного зеркала. Поверхностная зона ... 40 \par
18. Движение жидкости под уровнем грунтовых вод. Глубокие зоны ... 44 \par
19. Залегание погребенных вод 46 \par
20. Миграция погребенных вод 48 \par
21. Влияние уплотнения осадков на миграцию жидкостей 49 \par
22. Залегание газа и нефти 52 \par
23. Происхождение газа и нефти 52 \par
24. Миграция и скопление газа и нефти 54 \par
Глава П. Закон Дарси и измерение проницаемости пористой среды \par
1. Закон Дарси 58 \par
2. Степень обоснованности закона Дарси 59 \par
3. Постоянная в законе Дарси. Проницаемость пористой среды 70 \par
4. Единицы проницаемости и номенклатура 74 \par
5. Основы методики измерения и вычисления проницаемости пористой \par
среды 76 \par
6. Измерение проницаемости несцементированных песков 79 \par
7. Общая техника определений проницаемости сцементированной пористой \par
среды 83 \par
S. Дальнейшие подробности измерений проницаемости с помощью  \par
жидкостей 86 \par
9. Дальнейшие подробности измерений проницаемости с помощью газов 87 \par
10. Сравнение газа и жидкостей при измерениях проницаемости .... 88 \par
11. Измерения проницаемости в полевых условиях . 90 \par
12. Типовые значения проницаемости * 96 \par
13. Измерения пористости 100 \par
\smallskip
\smallskip
\smallskip
{[}Страница файла 625{]}\par
Оглавление \par
625 \par
Глава 111. Общие гидродинамические уравнения для течения \par
жидкостей в пористой среде \par
Стр. \par
1. Основные гидродинамические соотношения 107 \par
2. Классическая гидродинамика 110 \par
3. Обобщенная форма закона Дарси 112 \par
4. Уравнения движения 115 \par
5. Граничные и начальные условия 119 \par
6. Аналогия с остальными физическими проблемами 122 \par
7. Недекартовы системы координат 123 \par
8. Заключение 125 \par
часть вторая \par
УСТАНОВИВШЕЕСЯ ТЕЧЕНИЕ ЖИДКОСТЕЙ \par
Глава IV. Проблемы плоского течения и методы теории потенциала \par
1. Введение 128 \par
2. Радиальное течение в скважину 129 \par
3. Ряд Фурье 134 \par
4. Несимметричное течение в скважину 139 \par
5. Произвольное распределение давления на контуре 139 \par
6. Течение между неконцентричными круговыми границами. Функция \par
Грина 144 \par
7. Течение из бесконечного линейного источника питания в скважину. \par
Фронтальное продвижение. Метод отражений 149 \par
8. Течение из конечного линейного источника питания в песчаник  \par
бесконечной величины. Метод сопряженных функций 154 \par
9. Течение из конечного линейного источника питания в скважину.  \par
Преобразования сопряженной функции. Бесконечный ряд отображений . . 158 \par
10. Противодавление на плотину с расширенным основанием. Забивная \par
крепь отсутствует 163 \par
П. Противодавление на плотину с забивной крепью. Теорема Шварца-Кри- \par
стоффеля 165 \par
12. Фильтрационный расход под плотинами с удлиненной шириной  \par
основания. Отсутствие забивной крепи. Преобразование эллиптической  \par
функции , 175 \par
13. Фильтрационный расход под плотинами с удлиненной шириной  \par
основания при наличии забивной крепи 181 \par
14. Фильтрационный расход под коффердамом 185 \par
15. Анизотропная среда 189 \par
16. Выводы из общей теории потенциала. Теорема Грина 190 \par
17. Приближенные и неаналитические методы решения задач плоского  \par
течения 195 \par
18. Заключение 204 \par
Глава V. Пространственные (трехразмерные) задачи \par
1. Введение 216 \par
2. Сферическое течение 217 \par
3. Несовершенные скважины. Распределение потенциала 220 \par
4. Текущие дебиты несовершенных скважин 227 \par
5- Несовершенные скважины в анизотропных песчаниках 231 \par
6. Заключение 234 \par
Глава VI. Гравитационное течение \par
1. Введение 240 \par
2. Дренирование наклонного пласта песчаника. Задача Гопфа и Трефтца 244 \par
3. Решение задачи двухразмерного гравитационного течения методом  \par
годографа. Отображение границ 251 \par
4. Фильтрация черзз плотину с вертикальными фасами. Аналитическая \par
теория . . 2 53 \par
\smallskip
\smallskip
\smallskip
{[}Страница файла 626{]}\par
626 \par
Оглавление \par
Стр. \par
5. Численные приложения 258 \par
6. Изучение гравитационного течения с помощью электрических моделей 266 \par
7. Соответствие некоторых точных решений уравнения Лапласа для  \par
гравитационного течения 269 \par
8. Фильтрация воды из каналов и канав в песчаники с глубоко  \par
залегающим водяным зеркалом 276 \par
9. Фильтрация воды из каналов или канав в песчаники, которые  \par
подстилаются высокопроницаемыми гравийными ложами на малых глубинах 270 \par
10. Приближенная теория фильтрации воды через плотины с наклонными \par
фасами 281 \par
11. Фильтрационные струи из каналов и канав, поглощаемые неглубоко  \par
залегающим зеркалом воды 286 \par
12. Приближенное решение некоторых проблем ирригации и дренирования 292 \par
13. Субирригация 293 \par
14. Проблема водонасыщения 294 \par
15. Проблема эрозии 294 \par
16. Дренирование гончарными трубами 295 \par
17. Теория гравитационного течения Дюпюи-Форхгеймера 297 \par
18. Эксперименты на песчаных моделях с трехразмерными  \par
гравитационными течениями 302 \par
19. Составной напор при гравитационном течении 310 \par
20. Приближенная потенциальная теория расхода при гравитационном  \par
течении 312 \par
21. Заключение 318 \par
Глава VII. Системы с непостоянной проницаемостью \par
1. Введение. Поверхности разрыва непрерывности 331 \par
2. Непрерывное изменение величины проницаемости 332 \par
3. Прерывное радиальное изменение проницаемости 333 \par
4. Прилегающие слои с различной проницаемостью. Течение жидкости в \par
трещиноватых известняках 338 \par
5. Системы с ограниченными трещинами в известняках 345 \par
6. Теория солянокислотной обработки скважин из карбонатных  \par
коллекторов 347 \par
7. Эффект от солянокислотной обработки в радиальной системе 349 \par
8. Эффект от солянокислотной обработки высокотрещиноватых  \par
известняков 351 \par
9. Несовершенные скважины в переслаивающихся горизонтах .... 354 \par
10. Эффективность заиленного лайнера на эксплоатационную  \par
производительность скважины 362 \par
11. Заключение 369 \par
Глава VIII. Системы двух жидкостей \par
1. Введение 374 \par
2. Продвижение краевой воды. Общая природа и формулировка проблемы 374 \par
3. Линейное продвижение 380 \par
4. Двухразмерное радиальное продвижение 382 \par
5. Линия частиц жидкости в однородной системе 385 \par
6. Движение линейного контура в единичную скважину 387 \par
7. Прямое перемещение воды между двумя скважинами 390 \par
8. Влияние силы тяжести на геометрическую форму продвигающейся  \par
поверхности раздела 393 \par
9. Влияние разницы в величине вязкости между жидкостями с обеих  \par
сторон поверхности раздела 395 \par
10. Образование водяных конусов. Физические основы теории 397 \par
11. Аналитические выводы 402 \par
12. Подавление водяного конусообразования прослоями глин .... 409 \par
13. Заключение 411 \par
\smallskip
\smallskip
\smallskip
{[}Страница файла 627{]}\par
Оглавление \par
Глава IX. Многоскважинные системы \par
Стр. \par
1. Введение 418 \par
2. Малые группы скважин. Общая теория 420 \par
3. Примеры 421 \par
4. Зависимость эксплоатационной производительности от числа скважин \par
в группе 428 \par
5. Распределение давления на внешнем контуре 429 \par
6. Небольшие группы скважин, питающиеся бесконечным линейным  \par
контуром 431 \par
7. Бесконечные линейные ряды скважин 432 \par
8. Распределение давления у бесконечного ряда скважин. Напорная линия 433 \par
0. Двухлинейное расположение скважин. Эффект заслона 439 \par
10. Размещение скважин тремя рядами 446 \par
П. Размещение скважин в шахматном порядке 450 \par
12. Теория размещения отстоящих (внешних) скважин. Формулировка  \par
проблемы 451 \par
13. Однорядное размещение внешних скважин 452 \par
14. Многорядное размещение внешних скважин 456 \par
15. Числовой пример 460 \par
16. Проблема водной репрессии (флюдинг) 463 \par
17. Процесс образования поверхности раздела вода---нефть. Эксперименты \par
на электролитических моделях . . 464 \par
18. Эксперименты на модели с линейным контуром заводнения . . . 466 \par
19. Эксперименты, соответствующие искусственному методу водной  \par
репрессии 469 \par
20. Эффект барьеров в системе водной репрессии 473 \par
21. Модели с листовыми проводниками и распределение потенциала . . . 475 \par
22. Аналитические расчеты проводимости сеток скважин при водной  \par
репрессии. Общий метод 480 \par
23. Проводимость при прямолинейном заводнении 481 \par
24. Проводимость среды с пятью скважинами при водной репрессии 483 \par
25. Проводимость среды с семью скважинами при водной репрессии . . . 485 \par
26. Проводимость при шахматном размещении скважин 488 \par
27. Расчет коэфициента полезного действия водной репрессии 490 \par
28. Коэфициент полезного действия при водной репрессии с  \par
последовательным питанием от линейного контура 491 \par
29. Коэфициент полезного действия пятискважинного размещения . . 492 \par
30. Коэфициент полезного действия семискважинного размещения при  \par
водной репрессии : 493 \par
31. Коэфициент полезного действия шахматной расстановки при водной  \par
репрессии с питанием от линейного контура 494 \par
32. Общие наблюдения над сетками размещения при водной репрессии . . 495 \par
33. Сравнение сеток размещения при водной репрессии 496 \par
34. Заключение 502 \par
часть третья \par
НЕУСТАНОВИВШЕЕСЯ ТЕЧЕНИЕ ЖИДКОСТЕЙ \par
Глава X. Течение сжимаемых жидкостей в пористой среде \par
1. Введение 513 \par
2. Радиальное течение. Некоторые предварительные аналитические  \par
формулировки 520 \par
3. Радиальные системы, в которых плотности заданы на обоих границах 522 \par
4. Падение добычи в месторождении, которое эксплоатируется при  \par
гидравлическом режиме с переменным давлением 525 \par
5. Предельный случай нулевого внутреннего радиуса 529 \par
6. Рост забойных давлений в закрытых скважинах 529 \par
7. Радиальные системы с плотностью, заданной на одной границе и  \par
расходом на другой 531 \par
\smallskip
\smallskip
\smallskip
{[}Страница файла 628{]}\par
628 \par
Оглавление \par
Сет. \par
8. Падение давления в нефтяном месторожхеняи Ист-Тексас 532 \par
9. Единичная скважина в замкнутом резервуаре 54!» \par
10. Скважина с бесконечно малым радиусом ... 541 \par
11. Радиальные системы с заданными расходами на обоих границах . . . 541 \par
12. Предельный случай исчезающе малого внутреннего радиуса 54\textquotesingle Л \par
13. Скважина в замкнутом песчанике 544 \par
14. Нерадиал :Ное течение. Интерференция скважин. Функция Грина. . . 547 \par
15. Применение источников и стоков к решению проблем нестационарного \par
течения сжимаемых жидкостей в пористой среде 553 \par
16. Заключение 555 \par
ЧАСТЬ ЧЕТВЕРТАЯ \par
ДВИЖЕНИЕ ГАЗОВ В ПОРИСТОЙ СРЕДЕ \par
Глава XI. Движение газов в пористой среде \par
1. Введение 563 \par
2. Установившееся течение газов. Линейные системы 564 \par
3. Установившееся состояние двухразмерных (плоских) систем ... 565 \par
4. Трехразмерные (пространственные) системы в установившемся  \par
состоянии 566 \par
5. Влияние силы тяжести на течение газов в пористой среде 567 \par
6. Установившееся течение газов в системах с непостоянной  \par
проницаемостью 568 \par
7. Двухжидкостные системы. Водное конусообразование 569 \par
8. Газовые факторы в пористой среде, где газ и нефть движутся как  \par
однородные жидкости 570 \par
9. Газовые факторы для сообщающихся газовых и нефтяных зон . . 571 \par
10. Газовые факторы несообщающихся между собой газовых и нефтяных \par
песчаников 574 \par
11. Влияние спуска в скважину фонтанных трубок на величину газового \par
фактора 577 \par
12. Образование газовых конусов в скважинах с фонтанными трубками . . 581 \par
13. Многоскважинные системы 583 \par
14. Неустановившееся течение газов в пористой среде 584 \par
15. Замкнутый газоносный резервуар, дренируемый скважиной, которая \par
работает при постоянном давлении 587 \par
16. Замкнутый газоносный резервуар, дренируемый скважиной,  \par
работающей с постоянным отбором 591 \par
17. Заключение 592 \par
Приложение I. Уравнение Лапласа в криволинейных координатах . . . 599 \par
Приложение II. Некоторые двухразмерные функции Грина 600 \par
Приложение III. Преобразования с помощью модулярной эллиптической \par
функции 6(q) --- k*2 601 \par
Приложение IV. Доказательство обобщенной формулы Пуассона ... 602 \par
Приложение V. Сводка особых количественных выводов, развитых в  \par
настоящей работе в виде формул или графических построений ...... 603 \par
Список литературы русских исследователей по вопросу движения  \par
жидкости через пористую среду 616 \par
Основные пособия по математическому анализу (на русском языке). . . 623 \par
Редактор И. И. Гольберг Технический редактор А. С. Полосина \par
Подписано в печать 20.02.04. Формат 60 х 841/32- \par
Печать офсетная. Бумага офсетная № 1. \par
Усл. печ. л. 50,12. Уч. изд. л. 50,86. Гарнитура Тайме. \par
Научноиздательский центр «Регулярная и хаотическая динамика» \par
Отпечатано в полном соответствии с качеством \par
предоставленных диапозитивов в ГИПП «Вятка». \par
610033, г. Киров, ул. Московская, 122.\par
\endgroup

% ===== Раздел 12 =====

\section{Р. Коллинз. Течения жидкостей через пористые материалы. 1964}\label{section-12}

Файл: \texttt{Collins.djvu}.

\begingroup\small\setlength{\parindent}{0pt}\setlength{\parskip}{1pt}
{[}Страница файла 348{]}\par
ОГЛАВЛЕНИЕ \par
Глава 1 Структура и свойства пористых материалов 9 \par
1 10 Структура и классификация 9 \par
1 20 Структура и характеристики 11 \par
1 30 Пористость . 12 \par
1 31 Пористость и структура 12 \par
1 32 Влияние на пористость консолидации и уплотнения 15 \par
1 33 Измерение пористости . 16 \par
1 40 Удельная поверхность . 19 \par
1 41 Измерение удельной поверхности 20 \par
1 50 Проницаемость . 21 \par
1 51 Структурная интерпретация проницаемости 22 \par
1 52 Факторы, влияющие на проницаемость пористых \par
материалов 24 \par
1 60 Типичные значения пористости, проницаемости и удель \par
ной поверхности 25 \par
1 70 Распределение пористости и проницаемости в природных \par
пористых материалах 27 \par
1 80 Механические свойства пористых материалов 32 \par
Задачи 34 \par
Литература 35 \par
Глава 2 Статика жидкости в пористых средах . 36 \par
2 10 Насыщенность Жидкостью 36 \par
2 11 Измерение насыщенности 37 \par
2 20 Капиллярное давление . . 38 \par
2 21 Зависимость капиллярного давления от нчсыщенности 41 \par
2 22 Измерение капиллярного давления 42 \par
2 23 Капиллярный гистерезис 45 \par
2 24 Остаточная насыщенность (связанная вода) 47 \par
2 25 / функция Леверетта 48 \par
3 30 Распределение пор по размерам 49 \par
2 40 Давление паров жидкостей в пористых средах 50 \par
2 41 Определение зависимости давления пара от насыщен \par
ноети 54 \par
2 50 Теплота смачивания и теплота набухания 54 \par
2 60 Электрические характеристики насыщенных жидкостью \par
пористых сред 56 \par
347 \par
\smallskip
\smallskip
\smallskip
{[}Страница файла 349{]}\par
3 10 \par
3 11 \par
3 12 \par
3 13 \par
3 14 \par
3 15 \par
3 20 \par
3 30 \par
3 31 \par
3 32 \par
3 40 \par
3 50 \par
3 60 \par
2 70 Осмос . 60 \par
2 80 Электрохимические потенциалы 61 \par
2 90 Статическое распределение насыщенности в пористых \par
средах 62 \par
Задачи 63 \par
Литература 63 \par
Глава 3 Физическая и математическая теория течения \par
жидкостей в пористых средах 68 \par
Типы и механизмы течения жидкостей 68 \par
Ламинарное течение вязкой жидкости в пористой среде 68 \par
Турбулентное течение в пористых средах . 74 \par
Совместное ламинарное течение несмешивающихся жид \par
костей в пористой среде 76 \par
Фильтрационный потенциал и электроосмос 79 \par
Терчоосмос 82 \par
Ограничения классической гидродинамики 83 \par
Дифференциальная форма закона Дарси 83 \par
Совместное течение несмешивающихся жидкостей в по \par
ристой среде Относительные проницаемости 85 \par
Закон Дарси для анизотропных пористых сред 87 \par
Уравнения состояния жидкостей 91 \par
Уравнение неразрывности 92 \par
Дифференциальные уравнения течения жидкостей в по \par
ристых средах . 95 \par
3 70 Граничные условия 101 \par
Задачи 104 \par
Литература 104 \par
Глава 4 Стационарное ламинарное течение однородной \par
жидкости 105 \par
4 10 Установившийся режим 105 \par
4 20 Прямолинейный поток и измерение проницаемости 105 \par
4 30 Плоское горизонтальное течение в однородных средах \par
Математическая формулировка задачи 109 \par
4 31 Функция тока и компчексныи потенциал 113 \par
4 32 Радиальный поток между соосными круговыми цилинд \par
рами . . 117 \par
4 33 Конформное отображение 119 \par
4 34 Преобразование Кристоффеля --- Шварца Конформное \par
отображение прямоугольника 123 \par
4 40 Поперечная проницаемость цилиндрического керна 126 \par
4 50 Источники и стоки 131 \par
4 51 Система скважин вблизи места скачкообразного изме \par
нения характеристик пористой среды, метод изображе \par
ний 136 \par
4 60 Гравитационный дренаж свободные поверхности 140 \par
Задачи 142 \par
Литература . 143 \par
348 \par
\smallskip
\smallskip
\smallskip
{[}Страница файла 350{]}\par
Глава 5. Неустановившееся ламинарное течение  \par
однородной жидкости 144 \par
5.10. Неустановившееся течение сжимаемой жидкости . . 144 \par
5.20. Преобразование Лапласа и теорема Дюамеля 145 \par
5.30. Источники и стоки. Метод изображений 150 \par
5.40. Испытания по восстановлению давления в нефтяных \par
скважинах 155 \par
5.50. Поведение пластов при водонапорном режиме . . . , 164 \par
5.51. Бесконечный прямолинейный водоносный горизонт . . 167 \par
5.52. Бесконечный водоносный горизонт; радиальное течение 171 \par
5.53. Наклонный водоносный пласт 172 \par
5.60. Течени есовершенного газа. Численное интегрирование 176 \par
Задачи 182 \par
Литература 183 \par
Глава 6. Совместное течение несмешивающихся  \par
жидкостей 184 \par
6.10. Установившееся прямолинейное течение и измерение \par
относительных проницаемостей. Концевой эффект 184 \par
6.20. Вытеснение одной жидкости другой. Уравнение Бак- \par
ли --- Леверетта 189 \par
6.21. Интегрирование уравнений Бакли ---Леверетта  \par
методом Велджа. Вычисление относительных  \par
проницаемостей по данным о вытеснении 197 \par
6.30. Прямолинейное вытеснение с учетом влияния  \par
капиллярного давления 201 \par
6.40. Пропитка 213 \par
6.50. Многомерное двухфазное течение. Численные решения 220 \par
6.60. Турбулентное течение несмешивающихся жидкостей 221 \par
Задачи 222 \par
Литература 222 \par
Глава 7. Задачи с подвижными границами. Вытеснение. \par
Отложение твердых частиц 224 \par
7.10. Вытеснение одной жидкости другой 224 \par
7.11. Аналитическое решение задачи о прямолинейном  \par
вытеснении 228 \par
7.20. Двумерное вытеснение. Вытеснение водой при  \par
коэффициенте подвижности, равном единице 230 \par
7.21. Влияние коэффициента подвижности на величину  \par
коэффициента заводненности 246 \par
7.30. Фильтрация и отложение твердых частиц 249 \par
7.40. Неустойчивость движения фронта и образование языков 256 \par
Задачи 262 \par
Литература 262 \par
Глава 8. Совместное ламинарное течение смешивающихся \par
жидкостей 264 \par
8.10. Смешивающиеся и несмешивающиеся жидкости. Закон \par
диффузии Фика 264 \par
349 \par
\smallskip
\smallskip
\smallskip
{[}Страница файла 351{]}\par
8.20. Вытеснение из капиллярной трубки в случае, когда  \par
жидкости смешиваются 266 \par
8.30. Вытеснение в пористой среде. Случай смешивающихся \par
жидкостей 271 \par
8.31. Теория дисперсии в пористых средах 274 \par
8.40. Макроскопические характеристики вытеснения  \par
смешивающихся жидкостей 278 \par
8.50. Турбулентное течение смешивающихся жидкостей 281 \par
Задачи 281 \par
Литература 282 \par
Глава 9. Теория моделирования 283 \par
9.10. Понятие подобия 283 \par
9.20. Основы теории моделирования 285 \par
9.30. Масштабные модели и условия подобия 288 \par
9.40. Примеры применения анализа размерностей в случае, \par
когда известны уравнения, описывающие систему . . 290 \par
9.50. Пример аналогового моделирования 299 \par
9.60. Пример моделирования, когда уравнения,  \par
описывающие систему, неизвестны 303 \par
9.70. Частичное подобие в моделировании 306 \par
9.80. Специальные вопросы моделирования пористых сред 310 \par
9.90. Результаты изучения на моделях движения жидкостей \par
в пористых средах 311 \par
Приложение к главе 9 312 \par
Задачи 313 \par
Литература 313 \par
Глава 10. Течение с фазовыми переходами 315 \par
10.10. Течение с фазовыми переходами. Режим растворенного \par
газа 315 \par
10.20. Фазовое равновесие 315 \par
10.30. Пересыщение и недосыщение 325 \par
10.40. Математическое описание режима растворенного газа 329 \par
10.50. Другой подход к описанию режима растворенного газа 333 \par
10.60. Третий способ описания режима растворенного газа. \par
Уравнения материального баланса 336 \par
10.70. Общая схема описания течений с фазовыми переходами 340 \par
Задачи 341 \par
Литература 341 \par
Указатель 343\par
\endgroup

% ===== Раздел 13 =====

\section{Myron B. Allen. The Mathematics of Fluid Flow Through Porous Media. 2021}\label{section-13}

Файл: \texttt{Allen-The-mathematics-of-fluid-flow-through-porous.pdf}.

\begingroup\small\setlength{\parindent}{0pt}\setlength{\parskip}{1pt}
{[}Страница файла 9{]}\par
vii\par
Contents\par
Preface\par
xi\par
1\par
Introduction\par
1\par
1.1\par
Historical Setting\par
1\par
1.2\par
Partial Differential Equations (PDEs)\par
2\par
1.3\par
Dimensions and Units\par
4\par
1.4\par
Limitations in Scope\par
5\par
2\par
Mechanics\par
7\par
2.1\par
Kinematics of Simple Continua\par
7\par
2.1.1\par
Referential and Spatial Coordinates\par
7\par
2.1.2\par
Velocity and the Material Derivative\par
9\par
2.2\par
Balance Laws for Simple Continua\par
10\par
2.2.1\par
Mass Balance\par
11\par
2.2.2\par
Momentum Balance\par
13\par
2.3\par
Constitutive Relationships\par
16\par
2.3.1\par
Body Force\par
17\par
2.3.2\par
Stress in Fluids\par
17\par
2.3.3\par
The Navier--Stokes Equation\par
19\par
2.4\par
Two Classic Problems in Fluid Mechanics\par
20\par
2.4.1\par
Hagen--Poiseuille Flow\par
21\par
2.4.2\par
The Stokes Problem\par
23\par
2.5\par
Multiconstituent Continua\par
24\par
2.5.1\par
Constituents\par
25\par
2.5.2\par
Densities and Volume Fractions\par
26\par
2.5.3\par
Multiconstituent Mass Balance\par
29\par
2.5.4\par
Multiconstituent Momentum Balance\par
30\par
\smallskip
{[}Страница файла 10{]}\par
viii\par
Contents\par
3\par
Single-fluid Flow Equations\par
33\par
3.1\par
Darcy's Law\par
33\par
3.1.1\par
Fluid Momentum Balance\par
34\par
3.1.2\par
Constitutive Laws for the Fluid\par
35\par
3.1.3\par
Filtration Velocity\par
37\par
3.1.4\par
Permeability\par
38\par
3.2\par
Non-Darcy Flows\par
39\par
3.2.1\par
The Brinkman Law\par
39\par
3.2.2\par
The Forchheimer Equation\par
40\par
3.2.3\par
The Klinkenberg Effect\par
41\par
3.3\par
The Single-fluid Flow Equation\par
42\par
3.3.1\par
Fluid Compressibility and Storage\par
43\par
3.3.2\par
Combining Darcy's Law and the Mass Balance\par
44\par
3.4\par
Potential Form of the Flow Equation\par
44\par
3.4.1\par
Conditions for the Existence of a Potential\par
45\par
3.4.2\par
Calculating the Scalar Potential\par
46\par
3.4.3\par
Piezometric Head\par
47\par
3.4.4\par
Head-Based Flow Equation\par
48\par
3.4.5\par
Auxiliary Conditions for the Flow Equation\par
49\par
3.5\par
Areal Flow Equation\par
51\par
3.5.1\par
Vertically Averaged Mass Balance\par
52\par
3.5.2\par
Vertically Averaged Darcy's Law\par
54\par
3.6\par
Variational Forms for Steady Flow\par
55\par
3.6.1\par
Standard Variational Form\par
55\par
3.6.2\par
Mixed Variational Form\par
57\par
3.7\par
Flow in Anisotropic Porous Media\par
58\par
3.7.1\par
The Permeability Tensor\par
58\par
3.7.2\par
Matrix Representations of the Permeability Tensor\par
59\par
3.7.3\par
Isotropy and Homogeneity\par
61\par
3.7.4\par
Properties of the Permeability Tensor\par
62\par
3.7.5\par
Is Permeability Symmetric?\par
64\par
4\par
Single-fluid Flow Problems\par
67\par
4.1\par
Steady Areal Flows with Wells\par
67\par
4.1.1\par
The Dupuit--Thiem Model\par
67\par
4.1.2\par
Dirac 𝛿Models\par
70\par
4.1.3\par
Areal Flow in an Infinite Aquifer with One Well\par
73\par
4.2\par
The Theis Model for Transient Flows\par
75\par
4.2.1\par
Model Formulation\par
75\par
4.2.2\par
Dimensional Analysis of the Theis Model\par
76\par
4.2.3\par
The Theis Drawdown Solution\par
79\par
4.2.4\par
Solving the Theis Model via Similarity Methods\par
80\par
4.3\par
Boussinesq and Porous Medium Equations\par
84\par
\smallskip
{[}Страница файла 11{]}\par
Contents\par
ix\par
4.3.1\par
Derivation of the Boussinesq Equation\par
86\par
4.3.2\par
The Porous Medium Equation\par
88\par
4.3.3\par
A Model Problem with a Self-similar Solution\par
89\par
5\par
Solute Transport\par
95\par
5.1\par
The Transport Equation\par
95\par
5.1.1\par
Mass Balance of Miscible Species\par
96\par
5.1.2\par
Hydrodynamic Dispersion\par
97\par
5.2\par
One-Dimensional Advection\par
100\par
5.2.1\par
Pure Advection and the Method of Characteristics\par
101\par
5.2.2\par
Auxiliary Conditions for First-Order PDEs\par
103\par
5.2.3\par
Weak Solutions\par
104\par
5.3\par
The Advection--Diffusion Equation\par
106\par
5.3.1\par
The Moving Plume Problem\par
106\par
5.3.2\par
The Moving Front Problem\par
108\par
5.4\par
Transport with Adsorption\par
111\par
5.4.1\par
Mass Balance for Adsorbate\par
112\par
5.4.2\par
Linear Isotherms and Retardation\par
113\par
5.4.3\par
Concave-down Isotherms and Front Sharpening\par
114\par
5.4.4\par
The Rankine--Hugoniot Condition\par
116\par
6\par
Multifluid Flows\par
121\par
6.1\par
Capillarity\par
122\par
6.1.1\par
Physics of Curved Interfaces\par
122\par
6.1.2\par
Wettability\par
125\par
6.1.3\par
Capillarity at the Macroscale\par
127\par
6.2\par
Variably Saturated Flow\par
129\par
6.2.1\par
Pressure Head and Moisture Content\par
129\par
6.2.2\par
The Richards Equation\par
131\par
6.2.3\par
Alternative Forms of the Richards Equation\par
132\par
6.2.4\par
Wetting Fronts\par
133\par
6.3\par
Two-fluid Flows\par
134\par
6.3.1\par
The Muskat--Meres Model\par
134\par
6.3.2\par
Two-fluid Flow Equations\par
136\par
6.3.3\par
Classification of Simplified Flow Equations\par
137\par
6.4\par
The Buckley--Leverett Problem\par
139\par
6.4.1\par
The Saturation Equation\par
139\par
6.4.2\par
Welge Tangent Construction\par
141\par
6.4.3\par
Conservation Form\par
146\par
6.4.4\par
Analysis of Oil Recovery\par
146\par
6.5\par
Viscous Fingering\par
149\par
6.5.1\par
The Displacement Front and Its Perturbation\par
150\par
6.5.2\par
Dynamics of the Displacement Front\par
152\par
\smallskip
{[}Страница файла 12{]}\par
x\par
Contents\par
6.5.3\par
Stability of the Displacement Front\par
153\par
6.6\par
Three-fluid Flows\par
154\par
6.6.1\par
Flow Equations\par
156\par
6.6.2\par
Rock-fluid Properties\par
157\par
6.7\par
Three-fluid Fractional Flow Analysis\par
158\par
6.7.1\par
A Simplified Three-fluid System\par
159\par
6.7.2\par
Classification of the Three-fluid System\par
160\par
6.7.3\par
Saturation Velocities and Saturation Paths\par
162\par
6.7.4\par
An Example of Three-fluid Displacement\par
164\par
7\par
Flows With Mass Exchange\par
167\par
7.1\par
General Compositional Equations\par
168\par
7.1.1\par
Constituents, Species, and Phases\par
168\par
7.1.2\par
Mass Balance Equations\par
169\par
7.1.3\par
Species Flow Equations\par
170\par
7.2\par
Black-oil Models\par
172\par
7.2.1\par
Reservoir and Stock-tank Conditions\par
172\par
7.2.2\par
The Black-oil Equations\par
173\par
7.3\par
Compositional Flows in Porous Media\par
175\par
7.3.1\par
A Simplified Compositional Formulation\par
175\par
7.3.2\par
Conversion to Molar Variables\par
176\par
7.4\par
Fluid-phase Thermodynamics\par
178\par
7.4.1\par
Flash Calculations\par
178\par
7.4.2\par
Equation-of-state Methods\par
179\par
Appendix A\par
Dedicated Symbols\par
181\par
Appendix B\par
Useful Curvilinear Coordinates\par
183\par
B.1\par
Polar Coordinates\par
183\par
B.2\par
Cylindrical Coordinates\par
184\par
B.3\par
Spherical Coordinates\par
186\par
Appendix C\par
The Buckingham Pi Theorem\par
189\par
C.1\par
Physical Dimensions and Units\par
189\par
C.2\par
The Buckingham Theorem\par
190\par
Appendix D\par
Surface Integrals\par
193\par
D.1\par
Definition of a Surface Integral\par
193\par
D.2\par
The Stokes Theorem\par
194\par
D.3\par
A Corollary to the Stokes Theorem\par
195\par
Bibliography\par
197\par
Index\par
207\par
\endgroup

% ===== Раздел 14 =====

\section{Jacob Bear. Dynamics of Fluids in Porous Media. 1972}\label{section-14}

Файл: \texttt{Bear-Dynamics-of-fluids-in-porous-media-1972.djvu}.

\begingroup\small\setlength{\parindent}{0pt}\setlength{\parskip}{1pt}
{[}Страница файла 4{]}\par
Contents \par
Preface xvii \par
CHAPTER 1 \par
Introduction 1 \par
1.1 Aquifers, Ground Water and Oil Reservoirs 1 \par
1.1.1 Definitions 1 \par
1.1.2 The Moisture Distribution in a Vertical Profile 2 \par
1.1.3 Classification of Aquifers 4 \par
1.1.4 Properties of Aquifers 7 \par
1.1.5 The Oil Reservoir 8 \par
1.2 The Porous Medium 13 \par
1.3 The Continuum Approach to Porous Media 15 \par
1.3.1 The Molecular and Microscopic Levels 15 \par
1.3.2 Porosity and Representative Elementary Volume 19 \par
1.3.3 Areal and Linear Porosities 21 \par
1.3.4 Velocity and Specific Discharge 22 \par
1.3.5 Concluding Remarks 24 \par
CHAPTER 2 \par
Fluids and Porous Matrix Properties 27 \par
2.1 Fluid Density 27 \par
2.1.1 Definitions 27 \par
2.1.2 Mixture of Fluids 30 \par
2.1.3 Measurement of Density 31 \par
2.2 Fluid Viscosity 32 \par
2.2.1 Definition 32 \par
2.2.2 Non-Newtonian Fluids 33 \par
2.2.3 Units 34 \par
2.2.4 Effect of Pressure and Temperature 34 \par
2.2.5 Measurement of Viscosity 35 \par
2.3 Fluid Compressibility 37 \par
2.4 Statistical Description of Porous Media 38 \par
2.4.1 Particle-Size Distribution 39 \par
2.4.2 Pore-Size Distribution 41 \par
2.4.3 Other Statistical Descriptions 42 \par
\smallskip
\smallskip
\smallskip
{[}Страница файла 5{]}\par
2.5 Porosity 43 \par
2.5.1 Porosity and Effective Porosity 43 \par
2.5.2 Porosity, Structure and Packing 45 \par
2.5.3 Porosity Measurement 47 \par
2.6 Specific Surface 50 \par
2.6.1 Definitions 50 \par
2.6.2 Measurement of Specific Surface 51 \par
2.7 Matrix and Medium Compressibility 52 \par
CHAPTER 3 \par
Pressure and Piezometric Head 59 \par
3.1 Stress at a Point 59 \par
3.2 Hydrostatic Pressure Distribution 62 \par
3.3 Piezometric Head 63 \par
CHAPTER 4 \par
The Fundamental Fluid Transport Equations in Porous Media 65 \par
4.1 Particles, Velocities and Fluxes in a Fluid Continuum 65 \par
\textquotesingle{} 4.1.1 Definitions of Particles and Velocities 65 \par
4.1.2 Diffusive Velocities and Fluxes 68 \par
4.1.3 The Eulerian and Lagrangian Points of View 70 \par
4.1.4 The Substantial Derivative 7.1 \par
4.2 The General Conservation Principle 74 \par
4.3 Equations of Mass, Momentum and Energy Conservation in a Fluid \par
Continuum 77 \par
4.3.1 Mass Conservation of a Species 77 \par
4.3.2 Mass Conservation of a Fluid System 78 \par
4.3.3 Conservation of Linear Momentum of a Species a 79 \par
4.3.4 Conservation of Linear Momentum of a Fluid System 80 \par
4.4 Constitutive Assumptions and Coupled Processes 82 \par
4.4.1 General Considerations 82 \par
4.4.2 Principles to be Used in Forming Constitutive Equations 84 \par
4.4.3 Coupled Processes 85 \par
4.5 A Porous Medium Model 90 \par
4.5.1 The Conceptual Model Approach 90 \par
4.5.2 A Model of Flow Through a Porous Medium 92 \par
4.5.3 Frames of Reference 93 \par
4.5.4 An Averaging Procedure 95 \par
4.6 Equations of Volume and Mass Conservation 98 \par
4.6.1 Equation of Volume Conservation 98 \par
4.6.2 Equation of Mass Conservation of a Species in Solution 100 \par
4.6.3 Equation of Mass Conservation 102 \par
4.7 Equation of Motion 104 \par
\smallskip
\smallskip
\smallskip
{[}Страница файла 6{]}\par
4.8 Tortuosity and Permeability 106 \par
4.8.1 Relationship Between Tortuosity and Permeability 106 \par
4.8.2 Tortuosity and Other Transport Coefficients 112 \par
4.8.3 Formation Factor and Resistivity Index (Amyx 1960) in Reservoir \par
Engineering 113 \par
CHAPTER 5 \par
The Equation of Motion of a Homogeneous Fluid 119 \par
5.1 The Experimental Law of Darcy 119 \par
5.2 Generalization of Darcy\textquotesingle s Law 122 \par
5.2.1 Isotropic Medium 122 \par
5.2.2 Anisotropic Medium 123 \par
5.3 Deviations from Darcy\textquotesingle s Law 125 \par
5.3.1 The Upper Limit 125 \par
5.3.2 The Lower Limit 127 \par
5.3.3 The Slip Phenomenon 128 \par
5.4 Rotational and Irrotational Motion 129 \par
5.4.1 The Potential and Pseudopotential 129 \par
5.4.2 Irrotational Flow 131 \par
5.5 Hydraulic Conductivity of Isotropic Media 132 \par
5.5.1 Hydraulic Conductivity and Permeability 132 \par
5.5.2 Units and Examples 135 \par
5.6 Anisotropic Permeability 136 \par
5.6.1 The Principal Directions 137 \par
5.6.2 Directional Permeability 143 \par
5.7 Measurement of Hydraulic Conductivity 148 \par
5.7.1 General 148 \par
5.7.2 The Constant Head Permeameter 149 \par
5.7.3 The Falling Head Permeameter 150 \par
5.7.4 Determining Anisotropic Hydraulic Conductivity 150 \par
5.8 Layered Porous Media 151 \par
5.8.1 Flow Normal and Parallel to the Medium Layers 151 \par
5.8.2 Equivalent Hydraulic Conductivity of Arbitrarily Directed Flow . . 155 \par
5.8.3 A Layered Medium as an Equivalent Anisotropic Medium 156 \par
5.8.4 Girinskii\textquotesingle s Potential 157 \par
5.9 Compressible Fluids 159 \par
5.10 Derivations of Darcy\textquotesingle s Law 161 \par
5.10.1 Capillary Tube Models 162 \par
5.10.2 Fissure Models 164 \par
5.10.3 Hydraulic Radius Models 165 \par
5.10.4 Resistance to Flow Models 167 \par
5.10.5 Statistical Models 171 \par
5.10.6 Averaging the Navier-Stokes Equations 173 \par
5.10.7 Ferrandon\textquotesingle s Model 175 \par
\smallskip
\smallskip
\smallskip
{[}Страница файла 7{]}\par
5.11 Flow At Large Reynolds Numbers 176 \par
5.11.1 The Phenomenon 176 \par
5.11.2 Turbulence, Inertial Forces and Separation 177 \par
5.11.3 Some Examples of Proposed Nonlinear Motion Equations .... 182 \par
5.12 Seepage Forces and Stresses 184 \par
5.12.1 The Forces 184 \par
5.12.2 Piping and Quicksand 186 \par
CHAPTER 6 \par
Continuity and Conservation Equations for a Homogeneous Fluid ... 195 \par
6.1 The Control Volume 195 \par
6.2 Mass Conservation in a Nondeformable Porous Matrix 196 \par
6.2.1 The Basic Continuity Equation 196 \par
6.2.2 Continuity Equation for an Incompressible Fluid 198 \par
6.2.3 Continuity Equation for a Compressible Fluid 200 \par
6.3 Mass Conservation in a Consolidating Medium 202 \par
6.3.1 Vertical Compressibility Only 203 \par
6.3.2 Extension to Three Phases and to Three-Dimensional Consolidation . 208 \par
6.3.3 Barometric Efficiency of Aquifers 211 \par
6.4 Continuity Equations for Flow in Confined and Leaky Aquifers 213 \par
6.4.1 The Horizontal Flow Approximation 213 \par
6.4.2 Flow in a Confined Aquifer 214 \par
6.4.3 Flow in a Leaky Aquifer 216 \par
6.4.4 Averaging the Exact Equations over a Vertical Line 218 \par
6.4.5 The Boltzmann Transformation 221 \par
6.5 Stream Functions 222 \par
6.5.1 Pathlines, Streamlines, Streaklines and Fronts 223 \par
6.5.2 The Stream Function in Two-Dimensional Flow 225 \par
6.5.3 The Stream Functions in Three-Dimensional Flow 226 \par
6.5.4 The Partial Differential Equations for the Lagrange and Stokes Stream \par
Functions 230 \par
6.5.5 The Relationships between the Potential and the Stream Functions 233 \par
6.5.6 Solving Problems in the q\textgreater-ifi Plane 235 \par
6.6 Flow Nets and Ground Water Contour Maps 236 \par
6.6.1 The \textless p-tj) Flow Net 236 \par
6.6.2 The Ground Water Contour Map 239 \par
6.7 The Partial Differential Equations Describing Flow of an Inhomogeneous \par
Incompressible Fluid in Terms of W 241 \par
6.7.1 Two-Dimensional Flow 241 \par
6.7.2 Axisymmetric Flow 243 \par
CHAPTER 7 \par
Solving Boundary and Initial Value Problems 247 \par
7.1 Initial and Boundary Conditions 248 \par
\smallskip
\smallskip
\smallskip
{[}Страница файла 8{]}\par
7.1.1 Boundary of Prescribed Potential 250 \par
7.1.2 Boundary of Prescribed Flux 251 \par
7.1.3 The Steady Free (or Phreatic) Surface without Accretion 252 \par
7.1.4 The Unsteady Free (or Phreatic) Surface without Accretion .... 254 \par
7.1.5 The Steady Free (or Phreatic) Surface with Accretion 256 \par
7.1.6 The Unsteady Free (or Phreatic) Surface with Accretion 258 \par
7.1.7 Boundary of Saturated Zone (or of Capillary Fringe) 259 \par
7.1.8 The Seepage Face 260 \par
7.1.9 Capillary Exposed Faces 262 \par
7.1.10 Discontinuity in Permeability 263 \par
7.1.11 A Note on Anisotropic Media 269 \par
7.1.12 Boundary Conditions in Terms of Pressure or Density 270 \par
7.2 A Well Posed Problem 270 \par
7.3 Description of Boundaries in the Hodograph Plane 272 \par
7.3.1 The Hodograph Plane 272 \par
7.3.2 Boundaries in the Hodograph Plane 274 \par
7.3.3 Examples of Hodograph Representation of Boundaries 280 \par
7.3.4 Intersection of Boundaries of Different Types 284 \par
7.4 The Relations between Solutions of Flow Problems in Isotropic and  \par
Anisotropic Media 289 \par
7.4.1 The Flow Equations 290 \par
7.4.2 Relationships among Parameters in the Two Systems 291 \par
7.4.3 Examples 296 \par
7.5 Superposition and Duhamel\textquotesingle s Principles 297 \par
7.5.1 Superposition 297 \par
7.5.2 Unsteady Flow with Boundary Conditions Independent of Time . . 299 \par
7.5.3 Unsteady Flow with Time-Dependent Boundary Conditions .... 300 \par
7.6 Direct Integration in One-Dimensional Problems 301 \par
7.6.1 Solution of the One-Dimensional Continuity Equation 301 \par
7.6.2 Advance of a Wetting Front 303 \par
7.7 The Method of Images 304 \par
7.7.1 Principles 304 \par
7.7.2 Examples 306 \par
7.8 Methods Based on the Theory of Functions 312 \par
7.8.1 Complex Variables and Analytic Functions 313 \par
7.8.2 The Complex Potential and the Complex Specific Discharge 316 \par
7.8.3 Sources and Sinks 316 \par
7.8.4 Conformal Mapping 324 \par
7.8.5 The Schwarz-Christoffel Transformation 333 \par
7.8.6 Fictitious Flow in the w-Plane 337 \par
7.9 Numerical Methods 338 \par
7.9.1 Method of Finite Differences 338 \par
7.9.2 The Method of Finite Elements 346 \par
7.9.3 Relaxation Methods 348 \par
\smallskip
\smallskip
\smallskip
{[}Страница файла 9{]}\par
7.9.4 Schmidt\textquotesingle s Graphic Method 350 \par
7.10 Flow Nets by Graphic Methods 351 \par
CHAPTER 8 \par
Unconf ined Flow and the Dupuit Approximation 361 \par
8.1 The Dupuit Approximation 361 \par
8.1.1 The Dupuit Assumptions 361 \par
8.1.2 Examples of Application to Hydraulic Steady Flows in Homogeneous \par
Media 366 \par
8.1.3 Unconf ined Flow in an Aquifer with Horizontal Stratification . . . 369 \par
8.1.4 Unconf ined Flow in an Aquifer with Vertical Strata 372 \par
8.1.5 Unconf ined Flow in a Two-Dimensional Inhomogeneous Medium . . 373 \par
8.2 Continuity Equations Based on the Dupuit Approximation 374 \par
8.2.1 The Continuity Equation 374 \par
8.2.2 Boundary and Initial Conditions 378 \par
8.2.3 Some Solutions of Forchheimer\textquotesingle s Equation 379 \par
8.2.4 Some Solutions of Boussinesq\textquotesingle s Equation 381 \par
8.3 The Hodograph Method 388 \par
8.3.1 The Functions to and w 388 \par
8.3.2 The Hodograph Method 389 \par
8.3.3 Examples without a Seepage Face 391 \par
8.3.4 Hamel\textquotesingle s Mapping Function 398 \par
8.3.5 Zhukovski\textquotesingle s and Other Mapping Functions 403 \par
8.3.6 A Graphic Solution of the Hodograph Plane 406 \par
8.4 Linearization Techniques and Solutions 408 \par
8.4.1 First Method of Linearization of the Boussinesq Equation 408 \par
8.4.2 The Second Method of Linearization of the Boussinesq Equation . . 417 \par
8.4.3 The Third Method of Linearization of the Boussinesq Equation . . . 419 \par
8.4.4 The Method of Successive Steady States 420 \par
8.4.5 The Method of Small Perturbations 422 \par
8.4.6 The Shallow Flow Approximation 430 \par
CHAPTER 9 \par
Flow of Immiscible Fluids 439 \par
9.1 Introduction 439 \par
9.1.1 Types of Two-Fluid Flows 439 \par
9.1.2 The Abrupt Interface Approximation 439 \par
9.1.3 Occurrence 440 \par
9.2 Interfacial Tension and Capillary Pressure 441 \par
9.2.1 Saturation and Fluid Content 441 \par
9.2.2 Interfacial Tension and Wettability 441 \par
9.2.3 Capillary Pressure 444 \par
9.2.4 Drainage and Imbibition 449 \par
\smallskip
\smallskip
\smallskip
{[}Страница файла 10{]}\par
9.2.5 Saturation Discontinuity at a Medium Discontinuity 452 \par
9.2.6 Laboratory Measurement of Capillary Pressure 453 \par
9.3 Simultaneous Flow of Two Immiscible Fluids 457 \par
9.3.1 The Basic Motion Equations 457 \par
9.3.2 Relative Permeability 459 \par
9.3.3 Mass Conservation in Multiphase Flow 466 \par
9.3.4 Statement of the Multiphase Flow Problem 466 \par
9.3.5 The Buckley-Leverett Equations 468 \par
9.3.6 Simultaneous Flow of a Liquid and a Gas 472 \par
9.3.7 Laboratory Determination of Relative Permeability 473 \par
9.4 Unsaturated Flow 474 \par
9.4.1 Capillary Pressure and Retention Curve 475 \par
9.4.2 The Capillary Fringe 480 \par
9.4.3 Field Capacity and Specific Yield 483 \par
9.4.4 The Motion Equation 487 \par
9.4.5 Relative Permeability of Unsaturated Soils 491 \par
9.4.6 The Continuity Equation 495 \par
9.4.7 Methods of Solution and Examples 503 \par
9.4.8 Additional Comments on Infiltration and Redistribution of Moisture 513 \par
9.4.9 Comments on Vapor Movement in Unsaturated Flow 515 \par
9.5 Immiscible Displacement with an Abrupt Interface 519 \par
9.5.1 The Abrupt Interface Approximation 519 \par
9.5.2 Piezometric Heads and Dynamic Equilibrium Conditions at a Stationary \par
Interface 521 \par
9.5.3 The Boundary Conditions along an Interface 524 \par
9.5.4 Horizontal Interface Displacement 526 \par
9.5.5 Interface Displacement in the Vertical Plane 533 \par
9.5.6 Numerical and Graphic Methods 536 \par
9.5.7 Approximate Solutions based on Linearization 538 \par
9.5.8 Interface Stability 544 \par
9.6 Determining the Steady Interface by the Hodograph Method 547 \par
9.6.1 Boundary Conditions 548 \par
9.6.2 Description of Boundaries in the Hodograph Plane 549 \par
9.6.3 Examples 549 \par
9.7 The Interface in a Coastal Aquifer 557 \par
9.7.1 Occurrence 557 \par
9.7.2 The Ghyben-Herzberg Approximation 559 \par
9.7.3 Determining the Shape of a Stationary Interface by the Dupuit- \par
Ghyben-Herzberg Approximation 561 \par
9.7.4 Approximate Solution for the Moving Interface 563 \par
9.7.5 Interface Upconing 569 \par
9.7.6 The Dupuit-Ghyben-Herzberg Approximation for an Unsteady \par
Interface in a Thick Aquifer 573 \par
\smallskip
\smallskip
\smallskip
{[}Страница файла 11{]}\par
CHAPTER 10 \par
Hydrodynamic Dispersion 579 \par
10.1 Definition of Hydrodynamic Dispersion 579 \par
10.2 Occurrence of Dispersion Phenomena 582 \par
10.3 Review of Some Hydrodynamic Dispersion Theories 582 \par
10.3.1 Capillary Tube and Cell Models 583 \par
10.3.2 Statistical Models 587 \par
10.3.3 Spatial Averaging 603 \par
10.4 Parameters of Dispersion 605 \par
10.4.1 The Coefficients of Mechanical Dispersion and Hydrodynamic \par
Dispersion 605 \par
10.4.2 The Medium\textquotesingle s Dispersivity 611 \par
10.4.3 Dispersivity-Permeability Relationship 615 \par
10.5 The Governing Equations and Boundary Conditions 617 \par
10.5.1 The Partial Differential Equation in Cartesian Coordinates . . . .617 \par
10.5.2 The Partial Differential Equation in Curvilinear Coordinates . . . 619 \par
10.5.3 Initial and Boundary Conditions 622 \par
10.5.4 Solving the Boundary Value Problems 624 \par
10.5.5 The Use of Nondimensional Variables 626 \par
10.6 Some Solved Problems 626 \par
10.6.1 One-dimensional Flow 627 \par
10.6.2 Uniform Flow in a Plane 633 \par
10.6.3 Plane Radial Flow 634 \par
10.7 Heat and Mass Transfer 641 \par
10.7.1 Modes of Heat Transfer in a Porous Medium 641 \par
10.7.2 Formulation of the Problem of Heat and Mass Transfer in a Fluid \par
Continuum 643 \par
10.7.3 Formulation of the Problem of Heat and Mass Transfer in a Porous \par
Medium 644 \par
10.7.4 Comments on Some Heat and Mass Transfer Coefficients 647 \par
10.7.5 Simplifying the Macroscopic Heat and Mass Transfer Equations . 651 \par
10.7.6 Convective Currents and Instability 653 \par
10.7.7 Some Similitude Considerations 660 \par
CHAPTER 11 \par
Models and Analogs 665 \par
11.1 General 665 \par
11.2 Scaling Principles and Procedure 668 \par
11.2.1 The Two Systems 668 \par
11.2.2 Geometric Similarity 669 \par
11.2.3 Kinematic Similarity 670 \par
11.2.4 Dynamic Similarity 670 \par
11.2.5 Dimensional Analysis 671 \par
\smallskip
\smallskip
\smallskip
{[}Страница файла 12{]}\par
11.2.6 Inspectional Analysis 673 \par
11.2.7 Modified Inspectional Analysis 676 \par
11.3 The Sand Box Model 678 \par
11.3.1 Description 678 \par
11.3.2 Scales 680 \par
11.4 The Viscous Flow Analogs 687 \par
11.4.1 General 687 \par
11.4.2 Description of the Vertical Hele-Shaw Analog 687 \par
11.4.3 Establishing the Analogy between Analog and Prototype ..... 690 \par
11.4.4 Scales for the Vertical Analog 693 \par
11.4.5 Recommended Applications of Vertical Analog 696 \par
11.4.6 The Liquids 697 \par
11.4.7 The Horizontal Hele-Shaw Analog---Description and Scales . . . 697 \par
11.4.8\^{} Simulation of an Infinite Horizontal Aquifer 701 \par
11.5 Electric Analogs 702 \par
11.5.1 Description of the Electrolytic Tank and the Conducting Paper \par
Analogs 702 \par
11.5.2 Scales for the Electrolytic Tank Analog 708 \par
11.5.3 The Resistance Network Analog for Steady Flow 710 \par
11.5.4 The Resistance-Capacitance Network for Unsteady Flow 716 \par
11.5.5 The Ion Motion Analog 719 \par
11.6 The Membrane Analog 722 \par
11.7 Summary 725 \par
Answers to Exercises 729 \par
Bibliography 733 \par
Index 757\par
\endgroup

% ===== Раздел 15 =====

\section{Г. И. Баренблатт, В. М. Ентов, В. М. Рыжик. Движение жидкостей и газов в природных пластах. 1984}\label{section-15}

Файл: \texttt{Barenblatt-et-al-1984.pdf}.

\begingroup\small\setlength{\parindent}{0pt}\setlength{\parskip}{1pt}
{[}Страница файла 208{]}\par
ОГЛАВЛЕНИЕ\par
ПРЕДИСЛОВИЕ\par
Глава 1\par
ОСНОВНЫЕ ПОНЯТИЯ\par
ПОДЗЕМНОЙ \par
ГИДРО-\par
ГАЗОДИНАМИКИ\par
Глава 2\par
КЛАССИЧЕСКИЕ МО-\par
ДЕЛИ ТЕОРИИ ФИЛЬ-\par
ТРАЦИИ \par
ОДНОРОД-\par
НОЙ ж и д к о с т и\par
Глава 3\par
НЕКЛАССИЧЕСКИЕ\par
МОДЕЛИ ДВИЖЕНИЯ\par
ОДНОРОДНЫХ ЖИД-\par
КОСТЕЙ\par
Глава 4\par
ДВУХФАЗНАЯ ФИЛЬТ-\par
РАЦИЯ И ТЕОРИЯ ВЫ-\par
ТЕСНЕНИЯ НЕФТИ ВО-\par
ДОЙ\par
Глава 5\par
ФИЗИКО-ХИМИЧЕ-\par
СКАЯ \par
ПОДЗЕМНАЯ\par
ГИДРОДИНАМИКА\par
НЕФТЯНОГО \par
ПЛАСТА\par
СПИСОК ЛИТЕРАТУРЫ\par
§ 1. Особенности теории \par
движения \par
жидкости\par
и газа в природных пластах .\par
.\par
.\par
.\par
. \par
5\par
§ 2. Пористые \par
среды \par
\textquotesingle. \par
7\par
§ 3. Закон Дарси, пределы его применимости\par
и уточнения \par
9\par
§ 4. Уравнение \par
неразрывности \par
и \par
основные\par
уравнения теории фильтрации \par
14\par
§ 1. Простейшие установившиеся напорные те-\par
чения \par
. \par
20\par
§ 2. Качественные методы теории напорных те\^{}\par
чений \par
28\par
§ 3. Установившиеся \par
безнапорные \par
течения \par
36\par
§ 4. Нестационарное \par
движение \par
однородной\par
сжимаемой \par
жидкости. \par
Линейная \par
теория \par
38\par
§ 5. Нестационарное \par
движение \par
однородных\par
жидкостей. Нелинейные эффекты \par
.\par
.\par
. \par
51\par
§ 1. Теория фильтрации неньютоновских жид-\par
костей. Закон фильтрации \par
74\par
§ 2. Стационарные задачи фильтрации \par
ненью-\par
тоновских \par
жидкостей \par
.\par
.\par
.\par
.\par
.\par
.\par
. \par
79\par
§ 3. Нестационарные \par
задачи \par
фильтрации не-\par
ньютоновских \par
жидкостей \par
93\par
§ 4. Неравновесность \par
при фильтрации \par
одно-\par
родных \par
жидкостей. \par
Движение в трещи-\par
новато-пористых и \par
слоисто-неоднородных\par
пластах \par
104\par
§ 1. Основные представления теории двухфаз-\par
ного течения в пористых средах . . . . 1 1 8\par
§ 2. Структура двухфазного течения при круп-\par
номасштабном описании. Задача Баклея ---\par
Леверетта \par
128\par
§ 3. Структура течения при мелкомасштабном\par
описании. \par
Стабилизированная \par
зона. Ка-\par
пиллярные \par
эффекты в пористых \par
средах 136\par
§ 4. Неравновесные эффекты при двухфазной\par
фильтрации \par
150\par
§ 5. Устойчивость вытеснения несмешивающих-\par
ся жидкостей \par
155\par
§ 6. Теория вытеснения неньютоновских жид-\par
костей. \par
Влияние \par
вязко-пластических\par
свойств нефти на нефтеотдачу \textquotesingle{} \par
.\par
.\par
.\par
. 1 61\par
§ 1 Процессы \par
тепло- и массопереноса в по-\par
ристой \par
среде \par
172\par
§ 2. Вытеснение \par
нефти растворами \par
активных\par
примесей \par
. \par
182\par
§ 3. Эффекты диффузии и неравновесности в\par
задачах вытеснения нефти раствором ак-\par
тивной примеси \par
200\par
20"»\par
\endgroup

% ===== Раздел 16 =====

\section{M. Muskat. The Flow of Homogeneous Fluids Through Porous Media. Английское издание}\label{section-16}

Файл: \texttt{Muskat-eng.djvu}.

\begingroup\small\setlength{\parindent}{0pt}\setlength{\parskip}{1pt}
{[}Страница файла 12{]}\par
CONTENTS \par
Page \par
Preface vii \par
PART I. FOUNDATIONS \par
CHAPTER I \par
Introduction 1 \par
1.1. Scope of Subject -. 3 \par
1.2. Homogeneous Fluids 4 \par
1.3. Porous Media 6 \par
1.4. General Qualifications Governing the Analytical Theory. . 8 \par
1.5. The Nature of Porous Media 10 \par
1.6. Systematic Packing of Spheres 10 \par
1.7. The Packing of Natural Materials 13 \par
1.8. Compaction of Sand and Gravel 14 \par
1.9. Compaction in Clay 16 \par
1.10. Effect of Compaction and Alteration of Sediments on \par
Permeability 17 \par
1.11. Kinds of Rock and Their Fluid-bearing Properties .... 17 \par
1.12. Structure of Rocks in Relation to Underground Fluids . . 22 \par
1.13. Subsurface Fluids---Their Occurrence and Migratory \par
Behavior 26 \par
1.14. Occurrence of Ground Water 28 \par
1.15. Water Content of Unconsolidated Surface Sediments ... 30 \par
1.16. Fluid Movements in the Capillary Zone 30 \par
1.17. Fluid Movements below the Water Table. The Surface \par
Zone 34 \par
1.18. Fluid Movements below the Water Table. The Deep Zones 38 \par
1.19. The Occurrence of Connate Waters 40 \par
1.20. The Migration of Connate Waters 43 \par
1.21. The Influence of Compaction of Sediments on Migration of \par
Fluids. 44 \par
1.22. The Occurrence of Gas and Oil 47 \par
1.23. The Source of Petroleum---Gas and Oil 48 \par
1.24. The Migration and Accumulation of Gas and Oil 50 \par
CHAPTER II \par
Darcy\textquotesingle s Law and the Measurement of the Permeability of \par
Porous Media 55 \par
2.1. Darcy\textquotesingle s Law 55 \par
2.2. The Range of Validity of Darcy\textquotesingle s Law 56 \par
xiii \par
\smallskip
\smallskip
\smallskip
{[}Страница файла 13{]}\par
xiv \par
CONTENTS \par
P\textbackslash GK \par
2.3. The Constant in Darcy\textquotesingle s Law. The Permeability of a \par
Porous Medium 69 \par
2.4. Permeability Units and Nomenclature 74 \par
2.5. The Principle of the Method of Measurement and of the \par
Calculation of the Permeability of Porous Media.... 76 \par
2.6. The Measurement of the Permeability of Unconsolidated \par
Sands 79 \par
2.7. The General Technique for Permeability Determinations of \par
Consolidated Porous Media 85 \par
2.8. Further Details for Permeability Measurements with \par
Liquids 89 \par
2.9. Further Details for Permeability Measurements with Gases 91 \par
2.10. Gases versus Liquids for Permeability Measurements ... 91 \par
2.11. Permeability Measurements in the Field 93 \par
2.12. Typical Permeability Values 101 \par
2.13. Porosity Measurements 113 \par
CHAPTER III \par
General Hydrodynamical Equations for the Flow of Fluids \par
through Porous Media 121 \par
3.1. Fundamental Hydrodynamic Relations 121 \par
3.2. The Classical Hydrodynamics 125 \par
3.3. Generalized Form of Darcy\textquotesingle s Law 127 \par
3.4. The Equations of Motion 131 \par
3.5. Boundary and Initial Conditions 136 \par
3.6. Analogies with Other Physical Problems 139 \par
3.7. Other Coordinate Systems 141 \par
3.8. Summary 143 \par
PART II. STEADY-STATE FLOW OF LIQUIDS \par
CHAPTER IV \par
Two-dimensional Flow Problems and Potential-theory Methods 149 \par
4.1. Introduction 149 \par
4.2. The Radial Flow into a Well 150 \par
4.3. Fourier Series 156 \par
4.4. Unsymmetrical Flow into a Well 163 \par
4.5. Arbitrary Pressure Distributions over the Boundaries. . . 163 \par
4.6. The Flow between Nonconcentric Circular Boundaries. \par
Green\textquotesingle s Function 169 \par
4.7. The Flow from an Infinite Line Source into a Well. The \par
Line Drive. The Method of Images 175 \par
4.8. The Flow from a Finite Line Source into an Infinite Sand. \par
The Method of Conjugate Functions 181 \par
4.9. The Flow from a Finite Line Source into a Well. Conjugate- \par
function Transformations. Infinite Sets of Images . . . 186 \par
\smallskip
\smallskip
\smallskip
{[}Страница файла 14{]}\par
CONTENTS \par
xv \par
Page \par
4.10. The Uplift Pressure on a Dam of Extended Base Width. \par
No Sheet Piling 192 \par
4.11. The Uplift Pressure on a Dam with Sheet Piling Present. \par
The Schwarz Christoff el Theorem 195 \par
4.12. The Seepage Flux under Dams of Extended Base Width. \par
No Sheet Piling. Elliptic-function Transformations . . 208 \par
4.13. The Seepage Flux under Dams of Extended Base Width; \par
Sheet Piling Present 215 \par
4.14. The Seepage Flux underneath Coffer Dams 221 \par
4.15. Anisotropic Media 225 \par
4.16. General Potential-theory Results. Green\textquotesingle s Reciprocation \par
Theorem 227 \par
4.17. Approximate and Nonanalytic Methods of Solving Two- \par
dimensional Flow Problems 234 \par
4.18. Summary 244 \par
CHAPTER V \par
Theeb-dimensional Pkoblems 258 \par
5.1. Introduction 258 \par
5.2. Spherical Flow 259 \par
5.3. Partially Penetrating Wells. Potential Distributions. . . 263 \par
5.4. The Production Capacities of Partially Penetrating Wells . 271 \par
5.5. Partially Penetrating Wells in Anisotropic Sands 277 \par
5.6. Summary 281 \par
CHAPTER VI \par
Gbavity-flow Systems 287 \par
6.1. Introduction 287 \par
6.2. The Drainage of a Sloping Sand. The Problem of Hopf \par
and Trefftz 292 \par
6.3. The Treatment of Two-dimensional Gravity-flow Systems \par
by the Method of Hodographs. The Mapping of the \par
Boundaries 300 \par
6.4. The Seepage through a Dam with Vertical Faces; Analytical \par
Theory 303 \par
6.5. Numerical Applications 309 \par
6.6. The Study of Gravity-flow Systems by Means of Electrical \par
Models 318 \par
6.7. Some Exact Solutions of Laplace\textquotesingle s Equation Appropriate to \par
Gravity-flow Systems 323 \par
6.8. The Seepage of Water from Canals or Ditches into Sands \par
with Deep-lying Water Tables 331 \par
6.9. The Seepage of Water from Canals or Ditches into Sands \par
Underlain by Highly Permeable Gravel Beds at Shallow \par
Depths 336 \par
6.10. An Approximate Theory of the Seepage of Water through \par
Dams with Sloping Faces 338 \par
\smallskip
\smallskip
\smallskip
{[}Страница файла 15{]}\par
xvi CONTENTS \par
P\textbackslash GE \par
6.11. Seepage Streams from Canals and Ditches Which Merge \par
with Shallow Water Tables 346 \par
6.12. The Approximate Treatment of Some Irrigation and  \par
Drainage Problems 353 \par
6.13. Subirrigation ... 354 \par
6.14. The Problem of Water Logging 355 \par
6.15. The Erosion Problem 355 \par
6.16. Tile Drainage 356 \par
6.17. The Dupuit-Forchheimer Theory of Gravity-flow Systems. 359 \par
6.18. Sand-model Experiments with Three-dimensional Gravity- \par
flow Systems 365 \par
6.19. Composite Pressure-head and Gravity-flow Systems. . . . 375 \par
6.20. An Approximate Potential Theory of the Flux through \par
Gravity-flow Systems 377 \par
6.21. Summary 385 \par
CHAPTER VII \par
Systems op Nonuniform Permeability 400 \par
7.1. Introduction. Surfaces of Discontinuity 400 \par
7.2. Continuous Variations in the Permeability 401 \par
7.3. Discontinuous Radial Variations in the Permeability . . . 403 \par
7.4. Adjacent Beds of Different Permeability. Fluid Flow in \par
Fractured Limestones 409 \par
7.5. Bounded Limestone---Fracture Systems 417 \par
7.6. The Theory of the Acid Treatment of Limestone Wells . . 420 \par
7.7. The Effect of Acid Treatment in Radial Systems 422 \par
7.8. The Effect of Acid Treatment in Highly Fractured  \par
Limestones 425 \par
7.9. Partially Penetrating Wells in Stratified Horizons 429 \par
7.10. The Effect of a Sanded Liner on the Production Capacity \par
of a Well 439 \par
7.11. Summary 447 \par
CHAPTER VIII \par
Two-fluid Systems 453 \par
8.1. Introduction 453 \par
8.2. Edge-water Encroachment. General Nature and  \par
Formulation of the Problem 453 \par
8.3. Linear Encroachment 459 \par
8.4. Two-dimensional Radial P\textquotesingle ncroachment 402 \par
8.5. The History of a Line of Fluid Particles in a Homogeneous \par
System 466 \par
8.6. The Line Drive into a Single Well 468 \par
8.7. The Direct Drive between Two Wells 472 \par
8.8. The Effect of Gravity on the Shape of the Encroaching \par
Interface 476 \par
8.9. The Effect of the Differences in Viscosity between the \par
Fluids on the Two Sides of the Interface 478 \par
\smallskip
\smallskip
\smallskip
{[}Страница файла 16{]}\par
CONTENTS xvii \par
Page \par
8.10. Water Coning; Physical Basis of the Theory 480 \par
8.11. Analytical Development 487 \par
8.12. The Suppression of Water Coning by Shale Lenses .... 496 \par
8.13. Summary 499 \par
CHAPTER IX \par
Multiple-well Systems 507 \par
9.1. Introduction 507 \par
9.2. Small Groups of Wells. General Theory 509 \par
9.3. Examples 511 \par
9.4. The Dependence of the Production Capacity on the Number \par
of Wells in the Group 518 \par
9.5. The Pressure Distribution over the External Boundary . . 520 \par
9.6. Small Groups of Wells Supplied by an Infinite Line Drive. 521 \par
9.7. Infinite Sets of Wells in Linear Arrays 523 \par
9.8. The Pressure Distribution about an Infinite Array of Wells. \par
The Line Drive 524 \par
9.9. Two-line Arrays. Shielding Effects 530 \par
9.10. Three-line Arrays 540 \par
9.11. Staggered Wells in Line Arrays 545 \par
9.12. The\textquotesingle Theory of Offsetting. Statement of the Problem. . . 546 \par
9.13. Single-line Offsetting 548 \par
9.14. Multiple-line Offsetting 552 \par
9.15. Numerical Example 556 \par
9.16. The Problem of Water Flooding 560 \par
9.17. The Progress of the Water-oil Interface---Electrolytic- \par
model Experiments 561 \par
9.18. Model Experiments with Line Floods 565 \par
9.19. Experiments Corresponding to Artificial Flooding  \par
Operations 568 \par
9.20. Effect of Barriers in the Flooding System 572 \par
9.21. Conduction-sheet Models and Potential Distributions. . . 573 \par
9.22. Analytical Calculations of the Conductivities of Flooding \par
Network. General Method 580 \par
9.23. The Conductivity of the Direct Line-drive Flood 582 \par
9.24. The Conductivity of the Five-spot Flood 585 \par
9.25. The Conductivity of the Seven-spot Flood 587 \par
9.26. The Conductivity of the Staggered Line Drive 592 \par
9.27. The Calculation of Flooding Efficiencies 593 \par
9.28. The Efficiency of the Direct Line-drive Flood 594 \par
9.29. The Five-spot-flood Efficiency 596 \par
9.30. The Efficiency of the Seven-spot Flood 597 \par
9.31. The Efficiency of the Staggered Line-drive Flood 598 \par
9.32. General Observations on Flooding Networks 599 \par
9.33. Comparison of the Flooding Networks. . 600 \par
9.34. Summary 606 \par
\smallskip
\smallskip
\smallskip
{[}Страница файла 17{]}\par
xviii CONTENTS \par
PART III. THE NONSTEAD Y-STATE FLOW OF LIQUIDS \par
Page \par
CHAPTER X \par
The Flow of Compressible Liquids thkotjgh Porous Media , -. . 621 \par
10.1. Introduction 621 \par
10.2. Radial Flow Systems. Some Preliminary Analytical \par
Formulas \textquotesingle. . 630 \par
10.3. Radial Systems in Which the Density Is Specified over \par
Both Boundaries 632 \par
10.4. Production Decline in a Field Produced by a Water Drive \par
with Variable Field Pressures 635 \par
10.5. The Limiting Case of Vanishing Internal Radius 639 \par
10.6. The Rise of the Bottom-hole Pressures in Closed-in Wells . 640 \par
10.7. Radial Systems in Which the Density Is Specified over One \par
Boundary and the Flux over the Other 642 \par
10.8. The Pressure Decline in the "East Texas" Oil Field ... 644 \par
10.9. A Single Well in a Closed Reservoir 653 \par
10.10. A Well of Infinitesimal Radius 654 \par
10.11. Radial Systems in Which the Flux Is Specified over Both \par
Boundaries 655 \par
10.12. The Limiting Case of Vanishing Internal Radius 656 \par
10.13. A Well in a Closed Sand 657 \par
10.14. Nonradial Flow. Well Interference. Green\textquotesingle s Function. . 661 \par
10.15. The Use of Sources and Sinks in the Solution of Problems \par
Involving the Nonsteady-state Flow of Compressible \par
Liquids through Porous Media 666 \par
10.16. Summary 669 \par
PART IV. THE FLOW OF GASES THROUGH POROUS \par
MEDIA \par
CHAPTER XI \par
The Flow of Gases through Porous Media 679 \par
11.1. Introduction 679 \par
11.2. The Steady-state Flow of Gases. Linear Systems .... 680 \par
11.3. Two-dimensional Systems in the Steady State 681 \par
11.4. Three-dimensional Systems in the Steady State 683 \par
11.5. The Effect of Gravity on the Flow of Gases through Porous \par
Media 684 \par
11.6. The Steady-state Flow of Gases within Systems of  \par
Nonuniform Permeability 685 \par
11.7. Two-fluid Systems. Water Coning 686 \par
11.8. Gas-oil Ratios in Porous Media Carrying Gas and Oil as \par
Homogeneous Fluids 687 \par
11.9. Gas-oil Ratios for Communicating Gas and Oil Zones. . . 689 \par
\smallskip
\smallskip
\smallskip
{[}Страница файла 18{]}\par
CONTENTS \par
XIX \par
Pagk \par
11.10. Gas-oil Ratios for Non Communicating Gas and Oil Sands. 692 \par
11.11. The Effect of the Tubing on the Gas-oil Ratio 696 \par
11.12. Gas Coning in Tubed Wells 701 \par
11.13. Multiple-well Systems 704 \par
11.14. The Nonsteady-state Flow of Gases through Porous Media 704 \par
11.15. A Closed Gas Reservoir Drained by a Well Producing at \par
Constant Pressure 708 \par
11.16. A Closed Gas Reservoir Drained by a Well Producing at a \par
Uniform Rate 712 \par
11.17. Summary 714 \par
APPENDICES \par
I. Conversion Factors 723 \par
II. Laplace\textquotesingle s Equation in Curvilinear Coordinates 725 \par
III. Some Two-dimensional Green\textquotesingle s Functions 727 \par
IV. The Transformation Properties of the Modular Elliptic \par
Function 0(q) = k*2 729 \par
V. Proof of the Generalized Poisson Formula. Eq. 6.4(15) . . 732 \par
VI. Tabulation of the Specific Quantitative Results Developed \par
in This Work in Formula or Graphical Form 733 \par
Author Index 751 \par
Subject Index \par
755\par
\endgroup

% ===== Раздел 17 =====

\section{Рукописный конспект --- ориентир вместо оглавления}\label{section-17}

Файл: \texttt{porous-lecture-notes.pdf}. Отдельного оглавления нет. Ниже приблизительные тематические блоки по визуальному просмотру разворотов; границы требуют уточнения при конкретном вопросе. Рукописный текст плохо поддается OCR, поэтому не использовать его распознавание как точную цитату.

\begin{longtable}[]{@{}
  >{\raggedright\arraybackslash}p{(\columnwidth - 2\tabcolsep) * \real{0.5000}}
  >{\raggedright\arraybackslash}p{(\columnwidth - 2\tabcolsep) * \real{0.5000}}@{}}
\toprule\noalign{}
\begin{minipage}[b]{\linewidth}\raggedright
Страницы PDF
\end{minipage} & \begin{minipage}[b]{\linewidth}\raggedright
Тематический ориентир
\end{minipage} \\
\midrule\noalign{}
\endhead
\bottomrule\noalign{}
\endlastfoot
1--5 & Пористость, осреднение, скорость фильтрации, закон Дарси, балансовые уравнения \\
6--10 & Уравнения энергии, границы применимости закона Дарси, обобщения, соотношения на границах \\
11--14 & Несжимаемая жидкость, простые течения, приток к скважине, плоские задачи, комплексный потенциал \\
15--18 & Безнапорные течения, свободная поверхность, гидравлическая теория, фильтрация газа \\
19--22 & Упругий режим, волна давления, скважины; начало рассмотрения вытеснения \\
23--24 & Вытеснение и устойчивость границы раздела \\
25--30 & Двухфазная фильтрация, насыщенность, характеристики, разрывы, Бакли--Леверетт \\
31--33 & Капиллярность и структура переходной зоны, уравнения двухфазной фильтрации \\
\end{longtable}

\end{document}
